% BEGIN EMBEDDED BIBLIOGRAPHY
% Bibliography data are part of this source; the compiler writes its temporary input.
\begin{filecontents*}[overwrite]{\jobname.bib}
@unpublished{QED3Master,
 title={The field-strength {Osterwalder--Schrader} axioms for {QED} in three dimensions},
 date={2026-10-20}, note = {Comprehensive edition, 263 pp., 20 October 2026, before the modular integration revision}}
@unpublished{TerminalPressure,
 title={The terminal measure and pressure estimates for lattice {QED} in three dimensions},
 date={2026-10-20}, note={Companion manuscript, version of 20 October 2026}}
@unpublished{PositiveProbe,
 title={Positive compact probes and free response coefficients for lattice {QED} in three dimensions},
 date={2026-10-20}, note={Companion manuscript, version of 20 October 2026}}
@book{Hatcher2002,
 author={Hatcher, Allen}, title={Algebraic Topology}, year={2002},
 publisher={Cambridge University Press}}
@online{PostnikovLin2019,
 author={Postnikov, Alexander}, title={Algebraic Combinatorics},
 year={2019}, note={MIT 18.212, Lecture 23, notes by Andrew Lin, Proposition 4},
 url={https://ocw.mit.edu/courses/18-212-algebraic-combinatorics-spring-2019/f6113618a3e085ba86b799f65934886d_MIT18_212S19_lec23.pdf}}

@unpublished{ChamberLattice,
 title={Reflection positivity for chamber {QED} on the infinite lattice},
 date={2026-10-20}, note={Companion manuscript, 20 October 2026}}

@unpublished{CoverTuning,
 title={Agreement across covers and common tuning for lattice {QED} in three dimensions},
 date={2026-10-20}, note={Companion manuscript, 20 October 2026}}

@unpublished{FermionBlock,
 title={Chamber fermion locality and local integration},
 date={2026-10-20}, note={Companion manuscript, 20 October 2026}}

@unpublished{MaxwellBlock,
 title={Maxwell block minimization and Gaussian fluctuations},
 date={2026-10-20}, note={Companion manuscript, 20 October 2026}}

@unpublished{NonlinearStep,
 title={The exact nonlinear renormalization step for lattice {QED} in three dimensions},
 date={2026-10-20}, note={Companion manuscript, 20 October 2026}}

@unpublished{NormalizedContraction,
 title={Relevant extraction and contraction of local interactions},
 date={2026-10-20}, note={Companion manuscript, 20 October 2026}}

@unpublished{QED3Integration,
 title={The continuum field-strength construction for {QED} in three dimensions},
 date={2026-10-20}, note={Modular revision of 20 October 2026}}

@unpublished{VolumeTransfer,
 title={Volume transfer and component cumulants for lattice {QED} in three dimensions},
 date={2026-10-20}, note={Companion manuscript, 20 October 2026}}
\end{filecontents*}
% END EMBEDDED BIBLIOGRAPHY
% RRM1: exact sources and reference response.
\documentclass[11pt,letterpaper,reqno]{amsart}
% BEGIN INLINE manuscript-style
\makeatletter
\NeedsTeXFormat{LaTeX2e}[2020/10/01]
% Use with \documentclass[11pt,letterpaper,reqno]{amsart}.
% This file controls presentation only. It supplies no mathematical content.

\RequirePackage[T1]{fontenc}
\RequirePackage{lmodern}
\RequirePackage{microtype}
\RequirePackage[letterpaper,margin=1in,headheight=15pt]{geometry}
\RequirePackage{amsmath,amsthm,amssymb,mathtools}
\RequirePackage{enumitem}
\RequirePackage{perpage}
\RequirePackage{xcolor}
\definecolor{manuscriptred}{rgb}{0.6,0.1,0.1}
\definecolor{manuscriptblue}{rgb}{0.1,0.3,0.5}
\definecolor{manuscriptmagenta}{HTML}{60445C}
\RequirePackage[breaklinks=true,colorlinks=true,
  linkcolor=manuscriptred,citecolor=manuscriptblue,urlcolor=manuscriptmagenta]{hyperref}
\RequirePackage{cleveref}
\hypersetup{pdfauthor={}}

% Retain the AMS paragraph, title, heading and running-head conventions.
% Equations retain one global Arabic sequence, as in the reference.
\setcounter{tocdepth}{1}
\renewcommand{\thefigure}{\Roman{figure}}
\renewcommand{\thetable}{\Roman{table}}
\MakePerPage{footnote}
\renewcommand{\thefootnote}{%
  \ifnum\value{footnote}>9\arabic{footnote}\else\fnsymbol{footnote}\fi}

% Upright statements. Main results are A, B, C; supporting statements
% share a section-local sequence. These end marks indicate typography,
% never proof completeness, validation, or a deferred-proof status.
\theoremstyle{definition}
\newtheorem{ms@theorem}{Theorem}
\renewcommand{\thems@theorem}{\Alph{ms@theorem}}
\newtheorem*{ms@theoremstar}{Theorem}
\newtheorem{ms@proposition}{Proposition}[section]
\newtheorem{ms@lemma}[ms@proposition]{Lemma}
\newtheorem{ms@definition}[ms@proposition]{Definition}
\newtheorem{ms@corollary}[ms@proposition]{Corollary}
\newtheorem*{ms@corollarystar}{Corollary}

\newcommand{\ms@majorstart}{%
  \pushQED{\qed}\renewcommand{\qedsymbol}{$\blacksquare\square$}}
\newcommand{\ms@minorstart}{%
  \pushQED{\qed}\renewcommand{\qedsymbol}{$\blacksquare$}}
\newenvironment{theorem}[1][]{\ms@majorstart\ms@theorem[#1]}
  {\popQED\endms@theorem}
\newenvironment{theorem*}[1][]{\ms@majorstart\ms@theoremstar[#1]}
  {\popQED\endms@theoremstar}
\newenvironment{proposition}[1][]{\ms@minorstart\ms@proposition[#1]}
  {\popQED\endms@proposition}
\newenvironment{lemma}[1][]{\ms@minorstart\ms@lemma[#1]}
  {\popQED\endms@lemma}
\newenvironment{definition}[1][]{\ms@minorstart\ms@definition[#1]}
  {\popQED\endms@definition}
\newenvironment{corollary}[1][]{\ms@majorstart\ms@corollary[#1]}
  {\popQED\endms@corollary}
\newenvironment{corollary*}[1][]{\ms@majorstart\ms@corollarystar[#1]}
  {\popQED\endms@corollarystar}
% The source's lemma/proposition variants are deliberately not carried over:
% choose one marker for each statement type rather than encode proof status.

\theoremstyle{remark}
\newtheorem{remark}[ms@proposition]{Remark}
\newtheorem*{remark*}{Remark}
% The ordinary amsthm proof environment retains its open square.

\crefname{ms@theorem}{theorem}{theorems}
\Crefname{ms@theorem}{Theorem}{Theorems}
\crefname{ms@proposition}{proposition}{propositions}
\Crefname{ms@proposition}{Proposition}{Propositions}
\crefname{ms@lemma}{lemma}{lemmas}
\Crefname{ms@lemma}{Lemma}{Lemmas}
\crefname{ms@definition}{definition}{definitions}
\Crefname{ms@definition}{Definition}{Definitions}
\crefname{ms@corollary}{corollary}{corollaries}
\Crefname{ms@corollary}{Corollary}{Corollaries}
\crefname{remark}{remark}{remarks}
\Crefname{remark}{Remark}{Remarks}

% Bibliography configuration for this manuscript.
  \RequirePackage[backend=biber,style=alphabetic,backref=true,
    giveninits=true,isbn=false,date=year]{biblatex}
  \addbibresource{\jobname.bib}
  \setcounter{biburlnumpenalty}{100}
  \setcounter{biburllcpenalty}{100}
  \setcounter{biburlucpenalty}{100}
  \renewcommand*{\bibfont}{\footnotesize}
  \renewbibmacro{in:}{}
  \DeclareFieldFormat{pages}{#1}
  % Use the public alphabetic style. No private author-hash overrides;
  % repeated author names are printed in full.
\makeatother
% END INLINE manuscript-style
\usepackage{mathrsfs}
\newcommand{\R}{\mathbb R}
\newcommand{\C}{\mathbb C}
\newcommand{\Z}{\mathbb Z}
\newcommand{\E}{\mathbb E}
\newcommand{\dd}{\mathrm d}
\DeclareMathOperator{\supp}{supp}
\DeclareMathOperator{\tr}{tr}
\allowdisplaybreaks
\title[Source identities and reference response]{Exact source identities and normalized reference response for lattice QED in three dimensions}
\date{}
\hypersetup{pdftitle={Exact source identities and normalized reference response for lattice QED in three dimensions},pdfauthor={}}
\begin{document}
\begin{abstract}
We prove a cutoff-uniform estimate for the normalized external response
of three-dimensional lattice QED on its reference torus. The mass and
vacuum parameters are those of one common tuned trajectory. An exact
unfolding identifies the zero-flux law with a harmonic-averaged Gaussian
integral, and a source shift expresses Wick products as determinant
derivatives. We compare the terminal logarithm with its one-cell free
factors, estimate connected derivatives between separated cells, and
control the complex large-field complement. A finite fourth-order
cancellation then gives a positive power of the charge, uniformly in
the ultraviolet cutoff, on a nonempty interval of charges.
\end{abstract}
\maketitle
\tableofcontents
\section{Introduction}\label{sec:rrm-introduction}

We use the chamber regulator with weights $(6,6\sqrt2,6\sqrt2,6)$
and Wilson parameter $0<\rho<1$. The operator $D^{\rm ch}[B]$
acts on four-component spinors with $48$ chamber labels per site,
antiperiodic in every coordinate; its argument $B$ is the bond phase.
The stencil is recalled in \cref{subsec:rrm-regulator}.

Fix the block construction of \cite[Proposition~4.1 and \S4.1--4.4]{QED3Integration},
with its positive charge ceilings $e_1\le e_0<1$ and the
common-pair theorem \cite[Theorem~D]{CoverTuning}.
Thus $L\ge13$ is odd, $M=L^m$ with $m\in\Z_{\ge2}$,
$0<\alpha<1$, $1-\alpha<\varepsilon\le1/200$, and the localization and cutoff
powers are $r_{\rm loc}=2$, $p_0=9$, $p=19$.
The fixed ceilings denoted by $e_0,e_1$ here are the
$\bar e_0,\bar e_1$ of \cite{QED3Integration};
they are distinct from the running charge $e_j$ at $j=0$.
These fixed constants include the photon and fermion inverse bounds,
the exact block step and the terminal conditions. The parameter input is Proposition~4.1 of \cite{QED3Integration},
which supplies the construction and charge ceilings before any
probe-dependent restriction. The response estimates are proved here.

For $0<\bar m\le1$ put
\[
 h=L^{\lfloor\log_L\bar m^{-1}\rfloor},\qquad
 \eta=L^{-N},\qquad q=e\sqrt\eta,\qquad K=N+\log_Lh.
\]
Let $r_K$ be the least nonnegative integer power of $L$ at least
$[-\log(e\sqrt h)]^2$, and set
\[
 R_{\rm top}=hMr_K,\qquad R^*=R_{\rm top}/L,\qquad
 n_*=R^*/\eta,\qquad \ell_c=R^*/(2L),\qquad N_*=8L^3.
\]
The finite threshold $N_{\rm ref}(e)$ and stopping level on $R^*$
are those of \cite[Proposition~6.1 and (53)]{CoverTuning}.
For $N\ge N_{\rm ref}(e)$ use the real pair
$P_N(e)=(m_{0,N},\varepsilon_{0,N}^{\rm v})$ selected on
$R_{\rm top}$ by \cite[Theorem~D]{CoverTuning}, and write
$\boldsymbol\mu_N=\bar m\eta+m_{0,N}$.
This same pair is used for every source and every twist.
Its vacuum factor cancels from the normalized response below.

Fix a real nonzero $\beta\in C_c^\infty(\R^3;\Lambda^2)$,
put $f=dd^*\beta$, and let $D_\beta$ be the side of a least
coordinate cube containing $\supp\beta$. Enlarge that cube by two
in each direction to obtain $Q_\beta$ of side
$D_\beta^{\rm buf}=D_\beta+4$ and center $z_\beta$.
With $J_\eta$ denoting cube averaging and $\delta_{\eta,\mu}$
the centered difference, define the covering bond cochain
\begin{equation}\label{eq:rrm-probe-definition}
 \widehat a^o_\nu(x)=
 -\eta^{1/2}\sum_{\mu\ne\nu}
 (\delta_{\eta,\mu}J_\eta\beta_{\mu\nu})
                         (\eta x+\eta e_\nu/2).
\end{equation}
A placement $v\in\eta\Z^3$ is admitted when
$v+Q_\beta\subset[\ell_c/8,7\ell_c/8]^3$.
Translate the probe by $v$, reflect it into the cells indexed by
$\mathcal C_*=(\Z/2L\Z)^3$, and periodize on $C(n_*)$.
The resulting cochains are $\widehat a_c$; the reflection maps
and their orientation action are specified in \cref{sec:reference-cells}.
They have mutually orthogonal curvatures, all of squared norm
$v_\eta=\|d\widehat a^o\|_{\ell^2}^2$.

Write $C(n_*)=\Z^3/n_*\Z^3$ with its oriented cubical cells.
Let $\gamma_{n_*}$ be the standard Gaussian on its exact real
two-cochains, with counting inner product. For such a
curvature $u$, let $A(u)$ be its mean-zero transverse potential.
Set $H_{\theta,\nu}=\theta_\nu/(qn_*)$ and
$a(t)=e^{-1}\sum_ct_c\widehat a_c$. Define
\begin{equation}\label{eq:rrm-normalized-response}
 \mathring\Phi(t)=
 \frac{\displaystyle\int_{[-\pi,\pi]^3}\frac{\dd\theta}{(2\pi)^3}
       \int\det\bigl(D^{\rm ch}[q(A(u)+H_\theta+a(t))]
                         +\boldsymbol\mu_N\bigr)\,\dd\gamma_{n_*}(u)}
      {\displaystyle\int_{[-\pi,\pi]^3}\frac{\dd\theta}{(2\pi)^3}
       \int\det\bigl(D^{\rm ch}[q(A(u)+H_\theta)]
                         +\boldsymbol\mu_N\bigr)\,\dd\gamma_{n_*}(u)}.
\end{equation}
The denominator is positive, including when the tuned total mass
has either sign. The free one-cell ratio uses the positive reference mass:
\begin{equation}\label{eq:response-free-ratio}
 N_0(t)=\frac{\det\bigl(D^{\rm ch}[q(t/e)\widehat a_0]
                                      +\bar m\eta\bigr)}
                   {\det\bigl(D^{\rm ch}[0]+\bar m\eta\bigr)},
 \qquad c_m=N_0^{(m)}(0),\quad m=2,4.
\end{equation}
These are ordinary derivatives of the ratio. Put
\[
 \Theta_m^{(0)}=
 \left|\left[\prod_{c\in\mathcal C_*}
           (\partial_{t_c}^m-c_m)\right]\mathring\Phi(0)
                                      \right|^{1/N_*},\qquad
 \mathfrak e=\tfrac12(1/4-9\varepsilon)-12\varepsilon>0.
\]

\begin{theorem}[Normalized reference response]\label{thm:rrm-main}
For the fixed construction, mass $0<\bar m\le1$, and probe $\beta$
above, there is $e_{\rm RR}>0$ such that for every
$0<e\le e_{\rm RR}$ and every integer $N\ge N_{\rm ref}(e)$,
the function \eqref{eq:rrm-normalized-response} is entire and
\begin{equation}\label{eq:rrm-main-bound}
 \Theta_4^{(0)}\le C e^{\mathfrak e},\qquad
 (1+|c_2|)\Theta_2^{(0)}\le C e^{\mathfrak e}.
\end{equation}
The constants are independent of $N$ and of the admitted placement.
Every centered placement,
$|v+z_\beta-(\ell_c/2)(1,1,1)|_\infty\le\eta/2$, is admitted
in this range. One may take
$C=3\cdot10^5(1+\sqrt h\,C_\beta)^4h^{1/8}$,
where the finite, charge-independent $C_\beta$ is defined in
\eqref{eq:interacting-uniform-probe-constants}.

Let $\mathbb E^{(0)}$ denote the normalized zero-flux curvature law
on this same reference torus and at the same tuned mass.
Equivalently it is the probability law of $F^{\rm u}=u$ obtained
by weighting $\gamma_{n_*}$ with the zero-source twist-averaged
determinant in \eqref{eq:rrm-normalized-response}.
For $X_c=\langle d\widehat a_c,F^{\rm u}\rangle$ and
$H_2(y)=y^2-v_\eta$, $H_4(y)=y^4-6v_\eta y^2+3v_\eta^2$,
\begin{equation}\label{eq:rrm-main-wick}
 \mathbb E^{(0)}\prod_{c\in\mathcal C_*}
             (e^{-m}H_m(X_c)-c_m)
       =\left[\prod_{c\in\mathcal C_*}
           (\partial_{t_c}^m-c_m)\right]\mathring\Phi(0),
 \qquad m=2,4.
\end{equation}
All expectations are finite. The response estimate concerns the full
fixed family of $N_*=8L^3$ cells and the harmonic-averaged zero-flux law.
\end{theorem}

The proof has two distinct parts. The exact identities require no
renormalization estimates. Integer cochain unfolding retains the gauge
and harmonic Jacobians; a real source shift, followed by analytic
continuation, then gives \eqref{eq:rrm-main-wick}.
For the quantitative bound, the terminal density retains the external
connection throughout the common trajectory. Its logarithm is separated
into the free one-cell logarithms, their errors, and terms involving
several cells. A tree expansion preserves the distance between marked
cells and bounds the latter terms exponentially in their separation.

The restricted terminal density defines a positive measure of mass at
most one. The omitted configurations have small probability under the
zero-flux Gaussian bound. A two-constants estimate transfers this
smallness to the complex coefficient norm needed for differentiation.
In the ring retaining degree at most four in each source, every error
family which misses a cell is annihilated by that cell's free factor.
The surviving families therefore carry a power of the error for every
cell. This proves the estimate after taking the $N_*$th root.
The final section verifies all additional response inequalities on one
positive charge interval, including at jumps of the rounded scales.

These are the source and response arguments of
\cite[\S19--20, \S21.1, \S23]{QED3Master}.
The terminal operator estimates and exact Borel density used below are
proved in \cite[Theorem~B, \S6 and Appendix~A]{TerminalPressure}.
The free ratio has the convention of
\cite[Theorem~B]{PositiveProbe}; its coefficient estimates at the
analytic radius required here are proved in \cref{subsec:interacting-one-cell}.
No sign condition on the continuum fourth coefficient is needed.
Transfer to larger tori and other flux sectors is a subsequent result.

\subsection{Regulator and elementary conventions}\label{subsec:rrm-regulator}
Let $\sigma_\mu$ be the Pauli matrices and
$\gamma_\mu=\sigma_3\otimes\sigma_\mu$.
The labels $g\in O_{\rm cub}=O(3)\cap\mathrm{GL}(3,\Z)$ are
the signed permutation matrices. Let $s_0$ reverse coordinate three,
$s_1,s_2$ exchange $(2,3)$ and $(1,2)$, and let $r_{e_1}$ reverse
coordinate one. Put
\[
 n_0=e_3,\quad n_1=(e_2-e_3)/\sqrt2,\quad
 n_2=(e_1-e_2)/\sqrt2,\quad n_3=-e_1,\qquad
 M_i(g)=\tfrac12w_i(\rho+\gamma\cdot gn_i).
\]
On a unit torus the chamber operator is
\begin{equation}\label{eq:rrm-regulator}
 \begin{split}
 (D^{\rm ch}[B]\psi)(c,g)={}&
 \sum_{i=0}^2M_i(g)\bigl(\psi(c,g)-\psi(c,gs_i)\bigr)\\
 &+M_3(g)\bigl(\psi(c,g)
       -\mathrm e^{iB(c,c+ge_1)}\psi(c+ge_1,gr_{e_1})\bigr).
 \end{split}
\end{equation}
Reversed bonds carry the negative cochain value.
We write $D_A=D^{\rm ch}[qA]$ in the unit-lattice source identities.
At mesh $\eta_j$, the external hop and the bond phase have the
corresponding factors $\eta_j^{-1}$ and $e_j\eta_j$.
The spinor pairing has weight $1/48$ per chamber label; its constant
weight does not change a finite-dimensional determinant.

For smooth functions,
\[
 J_\eta v(y)=\eta^{-3}\int_{[-\eta/2,\eta/2]^3}v(y+z)\,\dd z,
 \qquad
 \delta_{\eta,\mu}v(y)=
 \frac{v(y+\eta e_\mu/2)-v(y-\eta e_\mu/2)}{\eta}.
\]
We use $\tau_vf(x)=f(x-v)$ and the seminorm
$|\beta|_{C^j}=\max_{\mu<\nu,|\alpha|\le j}
\|\partial^\alpha\beta_{\mu\nu}\|_\infty$.
Counting cochain pairings are bilinear after complexification.
Distances are periodic sup distances. The symbol $e$ is the physical
charge; numerical exponentials are written $\exp(\cdot)$ or $\mathrm e$.
\section{Probes on the reference cells}\label{sec:reference-cells}

The sampling and placement arguments are those of
\cite[\S19]{QED3Master}. The external-source identity uses a finitely supported bond cochain.
The continuum cumulant uses the test $f=dd^*\beta$. We identify the
discrete source with a nearby smooth test and place it strictly inside
one reference cell. Copies in distinct cells then have orthogonal
curvatures, which makes the normal-ordering factor separate over cells.

\subsection{The discrete source and the smooth test}

Fix a real nonzero $\beta\in C_c^\infty(\R^3;\Lambda^2)$.
Use the unbuffered support-cube side $D_\beta$ from the introduction,
and write
$D_\beta^{\mathrm{buf}}=D_\beta+4$. Choose a cube
$Q_\beta=z_\beta+[-D_\beta^{\mathrm{buf}}/2,
D_\beta^{\mathrm{buf}}/2]^3$ containing
$\supp\beta+[-2,2]^3$. Use the centered differences
$\delta_{\eta,\mu}$ and cube average $J_\eta$ defined in the introduction.
Put
\begin{align}
 a'_{\eta,\nu}&=-\sum_{\mu\ne\nu}
                          \delta_{\eta,\mu}\beta_{\mu\nu},\notag\\
 f^\eta_{\mu\nu}&=\delta_{\eta,\mu}a'_{\eta,\nu}
                         -\delta_{\eta,\nu}a'_{\eta,\mu}.
                                      \label{eq:cells-smooth-test}
\end{align}
All these functions are smooth and real. Since each centered
difference enlarges support by at most $\eta/2$ in its direction,
\[
 \supp f^\eta,\ \supp J_\eta a'_\eta
                   \subset\supp\beta+[-\eta,\eta]^3.
\]
In particular their supports lie in $Q_\beta$ for $0<\eta\le1$.

On the physical lattice $\eta\Z^3$ set
\[
 b_{\eta,\mu\nu}(c)
       =(J_\eta\beta_{\mu\nu})
                       (c+\eta(e_\mu+e_\nu)/2),\qquad
 a_\eta^o=d_\eta^*b_\eta.
\]
Here $d_\eta$ is the mesh-normalized forward cochain differential
and $d_\eta^*$ its adjoint for the pairing $\eta^3\sum$.
The adjoint formula gives
\begin{equation}\label{eq:cells-midpoint-source}
 a_{\eta,\nu}^o(c)=(J_\eta a'_{\eta,\nu})(c+\eta e_\nu/2).
\end{equation}
Transport to the unit lattice and define
\[
 \hat{\mathfrak a}_\eta^o(x,x+e_\nu)
                    =\eta^{1/2}a_{\eta,\nu}^o(\eta x).
\]
Its coboundary uses unscaled unit-lattice differences. The
midpoint identities yield
\begin{equation}\label{eq:cells-source-curvature}
 d\hat{\mathfrak a}_\eta^o(\sigma)
                        =f^\eta_{[\eta]}(\sigma).
\end{equation}
Indeed, the bond values in the curl are evaluated at the two
points displaced by $\pm\eta e_\mu/2$ from the plaquette center.
Their difference is $\eta\delta_{\eta,\mu}J_\eta a'_{\eta,\nu}$.
Commutation of the average and differences then gives
$\eta^{3/2}J_\eta f^\eta$ at that center, which is the right
side of \eqref{eq:cells-source-curvature}.

For a compactly supported smooth two-form $g$, let $g_{[\eta]}$
be the covering plaquette cochain with value $\eta^{3/2}J_\eta g$
at the plaquette center. On a torus $C(n)$ of physical side
$R=n\eta$, periodize this cochain by summing its translates by
$n\Z^3$. Write $F^\sharp$ for the periodic componentwise step
field whose value on the physical cube centered at a plaquette
is $\eta^{-3/2}F^{\rm u}$. Then $F^\sharp(g)$ is the finite-torus
pairing of the periodized $g_{[\eta]}$ with $F^{\rm u}$.
Periodizing the source likewise gives the exact observable identity
\begin{equation}\label{eq:cells-observable}
 \langle d\hat{\mathfrak a}_\eta,F^{\mathrm u}\rangle
                            =F^\sharp(f^\eta).
\end{equation}
The source is transverse: $a_\eta^o=d_\eta^*b_\eta$ is
orthogonal to every closed bond cochain, by summation by
parts. This applies to both the covering cochain and its
periodization.

The continuum approximation has the uniform bound
\begin{equation}\label{eq:cells-test-error}
 \|f^\eta-f\|_{H^s}\le\eta^2\|\beta\|_{H^{s+4}}
                         \qquad(s\in\R,\ 0<\eta\le1).
\end{equation}
To see this, the difference symbol is
$is_\mu(p)$ with $s_\mu(p)=2\eta^{-1}\sin(\eta p_\mu/2)$.
One has $|s_\mu|\le|p_\mu|$ and
$|s_\mu-p_\mu|\le\eta^2|p_\mu|^3/24$, so
\[
 |s_\mu s_\nu-p_\mu p_\nu|
                               \le\eta^2|p|^4/12.
\]
Each component of $f^\eta-f$ contains four such terms.
Cauchy--Schwarz on these terms, followed by summation over
the three two-form components, gives
$|\widehat{f^\eta-f}|^2\le(\eta^2/3)^2|p|^8|\hat\beta|^2$.
Plancherel's theorem proves \eqref{eq:cells-test-error}.

\subsection{Reference geometry and placement}

Let $T^*\ge1$ be the reference exponent furnished by the
construction, and put
\[
 R^*=hL^{T^*},\qquad n_*=R^*/\eta,\qquad
 p_{\rm cell}=n_*/L,\qquad w=p_{\rm cell}/2,\qquad \ell_c=\eta w=R^*/(2L).
\]
Since $h$ and $\eta^{-1}$ are powers of the odd integer $L$,
$p_{\rm cell}$ is an odd integer, $w\in\Z+1/2$, and the open unit-lattice
reference cells are $Q_c^\circ=w(c+(0,1)^3)$.
The reference torus $C(n_*)$ has $2L$ cells per side and
$N_*=8L^3$ in all.

The probe fits the reference cells if
\begin{equation}\label{eq:cells-fit}
 D_\beta^{\mathrm{buf}}+3\eta\le\ell_c.
\end{equation}
A placement is a vector $v\in\eta\Z^3$ satisfying
\[
 |v+z_\beta-(\ell_c/2)(1,1,1)|_\infty\le\eta/2.
\]
Such a vector always exists by coordinatewise rounding.
It translates $Q_\beta$ into $[\eta,\ell_c-\eta]^3$:
the displacement of its center is at most $\eta/2$, and
the half-side is at most $(\ell_c-3\eta)/2$.
Hence
\begin{equation}\label{eq:cells-probe-margin}
 \supp(\tau_v\beta)
                     \subset[2+\eta,\ell_c-2-\eta]^3.
\end{equation}
Form the source $\hat{\mathfrak a}^o$ with this translated
probe.

Every endpoint of a bond in the source support is within
$3\eta/2$ in sup distance of $\supp(\tau_v\beta)$.
By \eqref{eq:cells-probe-margin}, its lattice coordinates
therefore lie in
$[2/\eta-1/2,w-2/\eta+1/2]$.
If a plaquette has nonzero $d\hat{\mathfrak a}^o$, at least
one of its boundary bonds belongs to that support. Its
other two corners are one unit away. Thus the whole
plaquette lies in the box
\begin{equation}\label{eq:cells-support-box}
 B_0=(w/2)(1,1,1)+[-r_0,r_0]^3,\qquad
 r_0=w/2-2/\eta+3/2.
\end{equation}
Condition \eqref{eq:cells-fit} gives $r_0\ge3$.
The distance of this box from each bounding wall is
$2/\eta-3/2\ge1/2$, so $B_0\subset Q_0^\circ$.
Its reflected copies lie strictly inside the corresponding
cells. This is the support margin needed by the local
source identity.

\subsection{Orthogonal sources and translated observables}

For a cell index $c$, define the wall map $g_c$ coordinatewise:
if $c_\mu=2i$, it sends $x_\mu$ to $x_\mu+2iw$; if
$c_\mu=2i+1$, it sends $x_\mu$ to $(2i+2)w-x_\mu$.
It carries $Q_0^\circ$ to $Q_c^\circ$ and acts on cochains by
transporting oriented bonds, with a sign upon orientation reversal. Put
$\hat{\mathfrak a}_c^o=(g_c)_*\hat{\mathfrak a}^o$ on the
covering lattice and periodize it on $C(n)$, where $p_{\rm cell}$ divides
$n$ and $n\ge n_*$, to obtain $\hat{\mathfrak a}_c$.
The period lattice identifies cell indices modulo
$(2n/p_{\rm cell})\Z^3$. The coboundary commutes with these actions
and periodization, and its support lies in $B_c=g_cB_0$.
For cell representatives on this torus,
\begin{equation}\label{eq:cells-orthogonality}
 \langle d\hat{\mathfrak a}_c,d\hat{\mathfrak a}_{c'}\rangle_n
             =v_\eta\,\mathbf1_{\{c=c'\}},\qquad
 v_\eta=\|d\hat{\mathfrak a}^o\|^2.
\end{equation}
The covering norm is the unweighted $\ell^2$ norm, and the
torus pairing is the counting pairing.

To prove the identity, unfold one factor:
\[
 \langle d\hat{\mathfrak a}_c,d\hat{\mathfrak a}_{c'}\rangle_n
  =\sum_{\sigma}\sum_{z\in\Z^3}
       d\hat{\mathfrak a}_c^o(\sigma)
       d\hat{\mathfrak a}_{c'}^o(\sigma+nz).
\]
A nonzero summand has $\sigma\subset Q_c^\circ$ and
$\sigma+nz\subset Q_{c'}^\circ$. The open cells are disjoint,
so necessarily $c'=c+(2n/p_{\rm cell})z$. For distinct torus cells
there is no such $z$. For the same cell exactly one translate
can contribute; its summands are the squares of
$d\hat{\mathfrak a}_c^o$. The wall action preserves the
counting norm, which proves \eqref{eq:cells-orthogonality}.
Changing the placement by a vector in $\eta\Z^3$ translates
the covering source by an integer vector. Thus $v_\eta$
depends only on $\beta$ and $\eta$, and is independent of
the torus and the placement.

For a closed cochain $u$ define
$X_c(u)=\langle d\hat{\mathfrak a}_c,u\rangle_n$.
Every polynomial in this observable is read by one and the
same local integrand on all these tori. On the cover its
local expression is
\[
 \widehat A\longmapsto
 P\left(\sum_{\sigma\subset B_c}
             d\hat{\mathfrak a}_c^o(\sigma)(d\widehat A)(\sigma)\right).
\]
This is a scalar gauge-invariant polynomial. Cauchy--Schwarz
bounds its argument by
\[
 \sqrt{v_\eta}\left(\sum_{\sigma\subset B_c}
                      |d\widehat A(\sigma)|^2\right)^{1/2}.
\]
The polynomial is therefore bounded by
\[
 C_{\mathrm{pol}}\exp\left(\frac14\sum_{\sigma\subset B_c}
                                      |d\widehat A(\sigma)|^2\right)
\]
for a finite constant $C_{\mathrm{pol}}$ depending on the polynomial
and the covering source. This is the local Gaussian-growth condition used for the
observable class $\mathcal A(Q_c^\circ)$ in \cite[\S4.1, (46)]{ChamberLattice}.

The local boxes embed in the torus. Indeed, writing
$p_{\rm cell}=2m_{\rm cell}+1$, condition \eqref{eq:cells-fit} gives $m_{\rm cell}\ge7$,
and
\[
 2r_0+2=w-4/\eta+5\le w+1<p_{\rm cell}\le n.
\]
For each $c$, the translates of $B_c+[-1,1]^3$ by
$nz$, $z\in\Z^3$, are consequently disjoint.
On a local lift, $d\widehat A=F^{\mathrm u}$; unfolding
the finite cochain sum reads the polynomial as $P(X_c)$.
For a real polynomial its reflected copy from the base
cell is exactly $P(X_c)$, since conjugation fixes its
coefficients. This applies to the centered Wick polynomials.

The variance in \eqref{eq:cells-orthogonality} approaches
the continuum normal-ordering variance:
\begin{equation}\label{eq:cells-variance-limit}
 \left|\sqrt{v_\eta}-\|f\|_2\right|
 \le21\eta|\beta|_{C^3}(D_\beta^{\mathrm{buf}})^{3/2}
                      +\eta^2\|\beta\|_{H^4}.
\end{equation}
Let $\Pi_\eta f^\eta$ be the componentwise cell-average step
function on a period cube. By \eqref{eq:cells-source-curvature},
its $L^2$ norm is $\sqrt{v_\eta}$.
Each first derivative of $f^\eta_{\mu\nu}$ is bounded
by $4|\beta|_{C^3}$, from the four centered-difference
terms in \eqref{eq:cells-smooth-test} and the identity
$\delta_{\eta,\mu}g=\int_{-1/2}^{1/2}
(\partial_\mu g)(\,\cdot+t\eta e_\mu)\,\dd t$.
The pointwise error of the cell average is at most
$12\eta|\beta|_{C^3}$. It vanishes outside the periodic
images of the buffered cube, whose measure in one period
is $(D_\beta^{\mathrm{buf}})^3$.
Summing the three component squares bounds the $L^2$
error by $21\eta|\beta|_{C^3}(D_\beta^{\mathrm{buf}})^{3/2}$.
The norm triangle inequality and
\eqref{eq:cells-test-error} prove
\eqref{eq:cells-variance-limit}.

Finally the even cell indices give pure translations.
For $y\in(R^*/L)\Z^3=p_{\rm cell}\eta\Z^3$ set
$c(y)=2y/(p_{\rm cell}\eta)\in2\Z^3$. Then $g_{c(y)}$ is the unit-lattice
translation by $y/\eta$, and
\begin{equation}\label{eq:cells-translation-observable}
 X_{c(y)}=F^\sharp(\tau_{v+y}f^\eta).
\end{equation}
Two such cells agree on $C(n)$ exactly when their physical
translations agree modulo $R\Z^3$. Any fixed finite set
of distinct translations therefore occupies distinct
cells once the volume is sufficiently large.
Lattice translation invariance removes the common placement
$v$ from the expectation of any polynomial in these
observables. Its integrability follows from Gaussian
domination. These identities enter the subsequent volume transfer of
\cite[Theorem~A]{VolumeTransfer}, which supplies the larger-volume estimates
for the translates $\tau_yf^\eta$ and their passage to component cumulants.


\section{Exact external-response identities}\label{sec:external-response}

We prove the exact finite-volume identities of
\cite[\S20]{QED3Master}. We express the generating function of a curvature observable as a
determinant in an external connection. The identity is first proved
in each integer flux sector. In the zero-flux sector an unfolding
of the integer cochains replaces the compact photon integral by a
Gaussian integral over transverse cochains and an average over the
three twists. We compute its normalization before applying the
source shift.

\subsection{The determinant at complex connections}
\label{subsec:external-complex-determinant}

Use lattice units and fix $q>0$, $r_q=2\pi/q$ and a real lattice
mass $\mu_{\mathrm{lat}}$. Put
$M_{\mathrm{ch}}=\mu_{\mathrm{lat}}+\frac{\rho}{2}\sum_iw_i$.
For the law in \cref{thm:rrm-main},
$\mu_{\mathrm{lat}}=\boldsymbol\mu_N$.
The weights are the standard chamber weights, and $0<\rho<1$.
For an integer triple $n=(n_1,n_2,n_3)$ let
\[
 V=|n|=n_1n_2n_3,\qquad
 B_n=\prod_{\mu=1}^3\{0,\ldots,n_\mu-1\},\qquad
 \mathcal L_n=B_n\times\{1,2,3\}.
\]
In this subsection each $n_\mu$ may be as small as one.
Cochains are periodic functions on the oriented covering cells,
so the $3V$ link orbits remain distinct even at those small sides.
The chamber spinors are antiperiodic in each direction and have
dimension $192V$. Write $D_B$ for the massless chamber operator
with phases $\exp(iqB)$, and $\det_n$ for its finite-dimensional
determinant. In the phase convention of
\cref{subsec:rrm-regulator}, $D_B=D^{\mathrm{ch}}[qB]$.
The weight $1/48$ in the spinor pairing does not rescale this
determinant.

The determinant is a polynomial in the $6V$ variables
$\exp(iqB(b))$ and $\exp(-iqB(b))$, of degree at most $32$
in each. In particular it is entire in the complex link values.
To see the degree bound, choose the site basis on
$B_n\times O_{\mathrm{cub}}$ and include the antiperiodic sign
whenever a neighbor is reduced into $B_n$. Each of the four
rows at a site $(c,g)$ has constant entries except for the
one external-hop block. If $ge_1=e_\mu$, this block contains
$\exp(iqB_\mu(c))$; if $ge_1=-e_\mu$, it contains
$\exp(-iqB_\mu(c-e_\mu))$. For a fixed oriented photon link
there are eight labels at each end. Thus each phase variable
occurs in exactly $8\cdot4=32$ rows. The Leibniz determinant
chooses one entry in each row, proving the bound. The five
site labels of the diagonal and the four hops are distinct,
even if a side is one, so no other phase is hidden in one
of these blocks.

The two invariances needed below hold for complex cochains:
\begin{equation}\label{eq:external-determinant-invariances}
 \det_n(D_{B+r_qk}+\mu_{\mathrm{lat}})
       =\det_n(D_B+\mu_{\mathrm{lat}}),\qquad
 \det_n(D_{B+d\lambda}+\mu_{\mathrm{lat}})
       =\det_n(D_B+\mu_{\mathrm{lat}}),
\end{equation}
where $k$ is integer valued and $\lambda$ is a periodic complex
zero-cochain. The first identity follows directly from the phases.
For the second, multiplication by $\exp(iq\lambda(c))$ preserves
antiperiodicity and has inverse multiplication by
$\exp(-iq\lambda(c))$. On an external hop the factor
\[
 \mathrm e^{iq(B-d\lambda)(c,c')}\mathrm e^{iq\lambda(c')}
                                  =\mathrm e^{iq\lambda(c)}\mathrm e^{iqB(c,c')}
\]
gives $D_{B-d\lambda}=\mathrm e^{iq\lambda}D_B\mathrm e^{-iq\lambda}$.
Internal hops lie over one photon site and give the same
conjugation. This proves the determinant identity.

For $s\ge0$ define
\begin{align}
 K_n(s)&=
 \left(|M_{\mathrm{ch}}|+\frac{1+\rho}{2}
           \left(\sum_iw_i-6+6\mathrm e^{qs}\right)\right)^{192V},
           \notag\\
 |\det_n(D_B+\mu_{\mathrm{lat}})|
                   &\le K_n(\|\operatorname{Im}B\|_\infty).
                   \label{eq:external-determinant-strip}
\end{align}
Indeed each chamber bond involution permutes the site representatives.
For an internal hop its phase has modulus one; for an external
hop the modulus is at most $\mathrm e^{q\|\operatorname{Im}B\|_\infty}$.
The hopping matrix has norm $w_i(1+\rho)/2$. The operator norm
is therefore bounded by the base of the power in
\eqref{eq:external-determinant-strip}, and a determinant in dimension
$192V$ is bounded by that norm to the same power.
This estimate is uniform in the real part of the connection.

The twist connection is
\begin{equation}\label{eq:external-twist}
 H_\theta=\sum_\mu\frac{\theta_\mu}{qn_\mu}\mathbf1_\mu,
 \qquad
 \mathbf1_\mu=\sum_{c\in B_n}\mathbf1_{(c,\mu)}.
\end{equation}
Its determinant has period $2\pi$ in each complex twist variable.
For later use we need the cochain implementing this period. Set
\[
 k^{(\mu)}_\nu(c)=\delta_{\mu\nu}1_{\{c_\mu=n_\mu-1\}},
 \qquad
 \lambda^{(\mu)}(c)=-\frac{r_q}{n_\mu}c_\mu,\qquad c\in B_n,
\]
with periodic extension. On links not crossing the $\mu$ seam,
$d\lambda^{(\mu)}=-r_q/n_\mu$ in direction $\mu$ and zero
in the other directions. On a seam link its value is
$r_q-r_q/n_\mu$. Thus
\begin{equation}\label{eq:external-twist-period}
 H_{2\pi e_\mu}=r_qk^{(\mu)}-d\lambda^{(\mu)}.
\end{equation}
The two invariances in
\eqref{eq:external-determinant-invariances} prove the asserted
twist period. For $n_\mu=1$, $\lambda^{(\mu)}=0$ and
$k^{(\mu)}=\mathbf1_\mu$, so the same formula holds.

\subsection{Real and integer cochains on the torus}
\label{subsec:external-cochains}

From now on $n_\mu\ge4$. The counting pairings on cochains
are extended bilinearly when their arguments are complex.
Use
\[
 V_{\mathrm g}=dC^0(n),\qquad
 V_{\mathrm h}=\operatorname{span}_{\R}\{\mathbf1_1,\mathbf1_2,\mathbf1_3\},
 \qquad
 V_{\mathrm t}=(V_{\mathrm g}+V_{\mathrm h})^\perp.
\]
We recall the cochain arguments needed for the unfolding and
for its Jacobian. In particular the integer condition in the
flux sum requires an integer primitive, not merely a real one.

Let
\[
 \mathcal P_n^{\mathrm{per}}=
   \left\{p:p_\mu=2\pi j_\mu/n_\mu,\ 0\le j_\mu<n_\mu\right\},
 \qquad e_p(c)=V^{-1/2}\mathrm e^{ip\cdot c},\qquad
 \xi_\mu(p)=\mathrm e^{ip_\mu}-1.
\]
For the Fourier calculation we use the auxiliary Hermitian counting
inner product, under which the characters $e_p$ are orthonormal.
The analytic cochain pairings remain complex bilinear. At each momentum $p$ the
coboundary on exterior-valued cochains is exterior multiplication
by $\xi(p)$. If $\iota_{\bar\xi}$ denotes contraction by $\bar\xi$,
then
\[
 \iota_{\bar\xi}(\xi\wedge u)+
                  \xi\wedge\iota_{\bar\xi}u=|\xi|^2u.
\]
For $p\ne0$ a closed coefficient $u$ is consequently
$\xi\wedge(|\xi|^{-2}\iota_{\bar\xi}u)$, and is exact.
At zero momentum the cochains are constant and the
coboundary vanishes. Nonzero and zero modes are orthogonal,
and taking real parts of a complex primitive gives a real
primitive for a real cochain. Therefore
\begin{equation}\label{eq:external-Hodge}
 \ker(d:C^1\to C^2)=V_{\mathrm g}\oplus V_{\mathrm h},
 \qquad C^1=V_{\mathrm g}\oplus V_{\mathrm h}\oplus V_{\mathrm t},
 \qquad
 \ker(d:C^2\to C^3)=dC^1\oplus\operatorname{Harm}^2.
\end{equation}
Here $\operatorname{Harm}^2$ is the space of constant two-cochains.
The same Fourier formula gives $d^2=0$.

For a closed integer two-cochain $z$, define
$w_{\mu\nu}(z)$ as its sum over one coordinate $(\mu,\nu)$
two-torus. The sum is independent of the transverse coordinate:
sum $dz=0$ over the layer of cubes between two consecutive
planes. Differences in the two tangential directions cancel
by periodicity; the remaining difference is that of the two
plane sums. The fluxes give
\begin{equation}\label{eq:external-integer-two-exactness}
 z\in dC^1(n;\Z)\quad\Longleftrightarrow\quad
                         w_{\mu\nu}(z)=0\ \text{for all }\mu<\nu.
\end{equation}
The integer coefficients follow from the
cellular cohomology of the product of the three subdivided
circles. Its cochain complex is precisely the integer complex
used here. The three coordinate two-tori give a basis of
$H_2(T^3;\Z)$, and $H_1(T^3;\Z)$ is free. The universal
coefficient theorem identifies $H^2(T^3;\Z)$ with
$\operatorname{Hom}(H_2(T^3;\Z),\Z)$ by evaluation.
The class of $z$ evaluates to $w_{\mu\nu}(z)$ on each basis
two-torus, so it vanishes exactly when these sums vanish.
Vanishing of this cellular class is membership in
$dC^1(n;\Z)$.
\cite[Theorems~3.2, 3.5 and~3B.6]{Hatcher2002}

We next describe the closed integer one-cochains. If $d\lambda$
is integer valued, its sum along a path from $0$ to $c$ shows
that $\lambda(c)-\lambda(0)$ is integer valued. If $d\lambda=0$,
the same path argument shows that $\lambda$ is constant.
It follows that
\[
 C^1(n;\Z)\cap V_{\mathrm g}=dC^0(n;\Z).
\]
Each seam cochain $k^{(\mu)}$ in
\eqref{eq:external-twist-period} is closed, because its
multiple by $r_q$ lies in $V_{\mathrm g}+V_{\mathrm h}$.
For an arbitrary closed integer one-cochain $k$,
\eqref{eq:external-Hodge} gives
$k=d\lambda+\sum_\mu t_\mu\mathbf1_\mu$.
Its sum along a periodic $\mu$ loop is $n_\mu t_\mu$,
an integer $j_\mu$. Subtracting
$\sum_\mu j_\mu k^{(\mu)}$ leaves an integer exact cochain,
so
\begin{equation}\label{eq:external-closed-integer-one}
 k=d\lambda'+\sum_\mu j_\mu k^{(\mu)},\qquad
 \lambda'\in C^0(n;\Z),\quad j\in\Z^3.
\end{equation}
The harmonic projection of $k$ determines $j$, and then
determines $d\lambda'$. Moreover
\begin{equation}\label{eq:external-integer-Hodge}
 r_qk=
   \left(r_qd\lambda'+\sum_\mu j_\mu d\lambda^{(\mu)}\right)
                         +H_{2\pi j}.
\end{equation}
The first term is exact and the second is harmonic. Their
coupling through $j$ will matter when choosing a fundamental
domain.

The scalar Laplacian $\Delta_0=d^*d$ has eigenvalues
$|\xi(p)|^2=\sum_\mu(2-2\cos p_\mu)$ and a one-dimensional
kernel. Write
\begin{equation}\label{eq:external-scalar-determinant}
 \det{}'\Delta_0=\prod_{p\in\mathcal P_n^{\mathrm{per}}\setminus\{0\}}
                       |\xi(p)|^2,\qquad
 \mathcal L(n)=\log\det{}'\Delta_0.
\end{equation}
At each nonzero momentum the transverse one-cochains
are the two-dimensional space $\xi(p)^\perp$, with orthogonality
in the same Hermitian inner product. For $v\in\C^3$,
\[
 \sum_{\mu<\nu}|\xi_\mu v_\nu-\xi_\nu v_\mu|^2
           =|\xi|^2|v|^2-|\bar\xi\cdot v|^2.
\]
Thus $d^*d$ acts on that transverse space as multiplication
by $|\xi(p)|^2$, and
\begin{equation}\label{eq:external-transverse-determinant}
 \dim V_{\mathrm t}=2V-2,\qquad
 \det(d^*d|_{V_{\mathrm t}})=(\det{}'\Delta_0)^2.
\end{equation}
Diagonalizing this positive real form gives its Gaussian integral
\begin{equation}\label{eq:external-transverse-Gaussian}
 \mathcal G_n:=\int_{V_{\mathrm t}}\mathrm e^{-\|dA\|^2/2}\,\dd\sigma(A)
                    =\frac{(2\pi)^{V-1}}{\det{}'\Delta_0}.
\end{equation}
The measure $\sigma$ is Euclidean Lebesgue measure for the
counting pairing, not a coordinate measure chosen along a
nonorthogonal basis.

There is one more determinant needed for the exact directions.
The cochains $d\mathbf1_c$, $c\in B_n\setminus\{0\}$,
form both a real basis of $V_{\mathrm g}$ and an integer
basis of $dC^0(n;\Z)$. Indeed one subtracts the value at
zero from a zero-cochain; independence follows because
the resulting zero-cochain has value zero there and
cannot be a nonzero constant. Their Gram determinant is
\begin{equation}\label{eq:external-exact-Gram}
 \det\big(\langle d\mathbf1_c,d\mathbf1_{c'}\rangle\big)_{c,c'\ne0}
                           =\frac{\det{}'\Delta_0}{V}.
\end{equation}
To prove this formula, let $S$ be the scalar Laplacian matrix
in the site basis. The coefficient of $t$ in $\det(t+S)$
is $\det{}'\Delta_0$, by its eigenvalues. Multilinearity
of the determinant in columns identifies that coefficient
with the sum of the $V$ principal minors obtained by
deleting one row and column. All these minors have the same
determinant, since site translations conjugate $S$ by
permutations. The minor at zero is the Gram matrix above.

\subsection{Unfolding the zero-flux sector}
\label{subsec:external-zero-flux}

Let $\Phi$ be a Borel function on $C^2(n)$ with
$|\Phi(u)|\le C_\Phi \mathrm e^{C_\Phi\|u\|}$.
Let $W$ be a bounded continuous complex function on
$C^1(n)$, invariant under addition of real exact cochains
and of $r_q$ times integer one-cochains. Then
\begin{align}
 &\sum_{\substack{z\in C^2(n;\Z),\,dz=0\\w(z)=0}}
       \int_{[0,r_q)^{\mathcal L_n}}
            \Phi(dA+r_qz)\mathrm e^{-\|dA+r_qz\|^2/2}W(A)\,\dd A
       \notag\\
 &\quad=\mathcal N_n
       \int_{[-\pi,\pi]^3}\frac{\dd\theta}{(2\pi)^3}
       \int_{V_{\mathrm t}}
              \Phi(dA)\mathrm e^{-\|dA\|^2/2}W(A+H_\theta)\,\dd\sigma(A),
                                     \label{eq:external-zero-unfolding}\\
 &\mathcal N_n=r_q^{V+2}(\det{}'\Delta_0)^{1/2}.
                                     \label{eq:external-zero-normalization}
\end{align}
Both sides converge absolutely. The coefficient comes from
the lattice of exact gauge transformations and the three
periods of the harmonic directions; we give the unfolding
with both contributions.

First, set
\[
 C_q(\epsilon)=\sup_{y\in\R}\sum_{j\in\Z}
                  \mathrm e^{-\epsilon(y+r_qj)^2}
       \le2+r_q^{-1}\sqrt{\pi/\epsilon}\qquad(\epsilon>0).
\]
The bound follows by keeping the two nearest terms and comparing
the remaining terms with Gaussian integrals over adjacent intervals,
as in \cite[(5)]{ChamberLattice}. Summing each integer coordinate gives
\[
 \sum_{z\in C^2(n;\Z)}\int_{[0,r_q)^{\mathcal L_n}}
              \mathrm e^{-\epsilon\|dA+r_qz\|^2}\,\dd A
                     \le(r_qC_q(\epsilon))^{3V}.
\]
Here the integer closedness and flux restrictions have
been dropped. Completing the scalar square bounds
$|\Phi(u)|\mathrm e^{-\|u\|^2/2}$ by
$C_\Phi \mathrm e^{C_\Phi^2}\mathrm e^{-\|u\|^2/4}$.
Together with boundedness of $W$, this proves absolute
convergence of the left side.

Let $\mathcal Z^1_{\Z}=\ker d\cap C^1(n;\Z)$ and
choose representatives $\mathcal R$ of
$C^1(n;\Z)/\mathcal Z^1_{\Z}$. By
\eqref{eq:external-integer-two-exactness}, $k\mapsto dk$
is a bijection from $\mathcal R$ onto the closed
zero-flux integer two-cochains. Put
\[
 P(A)=\Phi(dA)\mathrm e^{-\|dA\|^2/2}W(A),\qquad
 \mathcal D_1=\bigcup_{k\in\mathcal R}
                     ([0,r_q)^{\mathcal L_n}+r_qk).
\]
The two invariances of $W$ turn the left side of
\eqref{eq:external-zero-unfolding} into
$\int_{\mathcal D_1}P(A)\,\dd A$.
The boxes are disjoint. Every $A$ has a unique decomposition
into its coordinatewise remainder in $[0,r_q)$ and
$r_q$ times an integer cochain. Decomposing that cochain
into its chosen representative and a closed one shows
that $\mathcal D_1$ is a Borel fundamental domain for
$r_q\mathcal Z^1_{\Z}$. The function $P$ is invariant
under that group.

We replace this domain by
\[
 \mathcal F_{\mathrm g}
 =\left\{r_q\sum_{c\ne0}t_c\,d\mathbf1_c:0\le t_c<1\right\},
 \qquad
 \mathcal F_{\mathrm h}=\{H_\theta:\theta\in[-\pi,\pi)^3\},
 \qquad
 \mathcal D_2=\mathcal F_{\mathrm g}+\mathcal F_{\mathrm h}
                                      +V_{\mathrm t}.
\]
To check that it is a fundamental domain for the same
group, write $A=A_{\mathrm g}+A_{\mathrm h}+A_{\mathrm t}$.
Formula \eqref{eq:external-integer-Hodge} shows that
subtracting $r_qk$ changes the harmonic part by
$H_{2\pi j}$ and the exact part by
$\sum_\mu j_\mu d\lambda^{(\mu)}+r_qd\lambda'$.
There is exactly one $j$ placing the harmonic part
in $\mathcal F_{\mathrm h}$. For this $j$ there is
exactly one $d\lambda'$ placing the exact part
in $\mathcal F_{\mathrm g}$, by the integer basis
used in \eqref{eq:external-exact-Gram}. The transverse
part is unchanged. Uniqueness in
\eqref{eq:external-closed-integer-one} gives uniqueness
of $k$. The orthogonal decomposition identifies
$\mathcal D_2$ with a product of Borel sets, so it
is Borel.

For an integrable periodic Borel function, the integral
over two such fundamental domains is the same.
Indeed partition $\mathcal D_1$ into its intersections
with $\mathcal D_2+h$, $h\in r_q\mathcal Z^1_{\Z}$,
and translate each intersection by $-h$. They partition
$\mathcal D_2$, and invariance of both the measure and
the function gives equality. Tonelli proves it first
for nonnegative functions; applying it to the modulus
then justifies the complex case by absolute integrability.
Thus $\int_{\mathcal D_1}P=\int_{\mathcal D_2}P$.

The two finite-dimensional volumes are
\begin{align}
 \operatorname{vol}(\mathcal F_{\mathrm g})
    &=r_q^{V-1}(\det{}'\Delta_0/V)^{1/2},
                                     \label{eq:external-gauge-volume}\\
 \dd A_{\mathrm h}
    &=V^{1/2}q^{-3}\,\dd\theta
      =r_q^3V^{1/2}\frac{\dd\theta}{(2\pi)^3}.
                                     \label{eq:external-twist-Jacobian}
\end{align}
The first follows from the square root of the Gram
determinant. For the second, the constant cochains
$V^{-1/2}\mathbf1_\mu$ are orthonormal, and the
matrix of $\theta\mapsto H_\theta$ in this basis
has diagonal entries $V^{1/2}/(qn_\mu)$.
Their product is $V^{1/2}q^{-3}$.
Their combined factor is precisely $\mathcal N_n$.
Since $P(A_{\mathrm g}+A_{\mathrm h}+A_{\mathrm t})
=\Phi(dA_{\mathrm t})\mathrm e^{-\|dA_{\mathrm t}\|^2/2}
W(A_{\mathrm t}+A_{\mathrm h})$, it is independent
of $A_{\mathrm g}$. Fubini over the orthogonal
decomposition and the two displayed volume formulas
prove \eqref{eq:external-zero-unfolding}, including
absolute convergence of its right side.

Apply this identity to
$W(A)=\det_n(D_A+\mu_{\mathrm{lat}})$. Its hypotheses
follow from \eqref{eq:external-determinant-invariances}
and \eqref{eq:external-determinant-strip}. If
$Z_n^{(0)}[\Phi]$ is the zero-flux part of the
configuration integral, we have proved
\begin{equation}\label{eq:external-zero-sector}
 Z_n^{(0)}[\Phi]=\mathcal N_n
       \int\frac{\dd\theta}{(2\pi)^3}
       \int_{V_{\mathrm t}}\Phi(dA)\mathrm e^{-\|dA\|^2/2}
          \det_n(D_{A+H_\theta}+\mu_{\mathrm{lat}})\,\dd\sigma(A).
\end{equation}
The twist integral always runs over $[-\pi,\pi]^3$.
For real $A$ and $\theta$ the determinant is
nonnegative. Also $Z_n^{(0)}[1]>0$ at every real
$\mu_{\mathrm{lat}}$: the lower bound in
\cite[Proposition~2.3]{ChamberLattice} was obtained
from the single term $z=0$, which belongs to this
sector. Dividing \eqref{eq:external-zero-sector}
by its value at $\Phi=1$ identifies the zero-flux
curvature law with the law of $dA$ under the
harmonic-averaged transverse measure.

\subsection{The Gaussian normalization and matter density}
\label{subsec:external-Gaussian-law}

Let $F_n=dC^1(n)=dV_{\mathrm t}$ and let $\gamma_n$
be the standard Gaussian probability measure on this
Euclidean space. The map $d_{\mathrm t}=d|_{V_{\mathrm t}}$
is a linear bijection onto $F_n$, by
\eqref{eq:external-Hodge}; write $A(u)=d_{\mathrm t}^{-1}u$.
Changing variables along this map gives
\begin{equation}\label{eq:external-Gaussian-change}
 \int_{V_{\mathrm t}}G(dA)\mathrm e^{-\|dA\|^2/2}\,\dd\sigma(A)
                     =\mathcal G_n\int G(u)\,\dd\gamma_n(u).
\end{equation}
Indeed the Jacobian is $|\det d_{\mathrm t}|^{-1}$,
and $\dim F_n=2V-2$. The Gaussian density contributes
$(2\pi)^{V-1}$; setting $G=1$ identifies the constant
with \eqref{eq:external-transverse-Gaussian}.
The formula holds for nonnegative Borel functions,
and for the exponentially bounded complex functions
used above by their Gaussian integrability.

Define the photon factor and normalized matter factor by
\begin{align}
 Z_n^{\mathrm{ph}}&=\mathcal N_n\mathcal G_n
       =r_q^{V+2}(2\pi)^{V-1}
                            \mathrm e^{-\mathcal L(n)/2},
                                  \label{eq:external-photon-factor}\\
 z_{r,n}(q,\mu_{\mathrm{lat}})
   &=\int\frac{\dd\theta}{(2\pi)^3}
       \int\det_n(D_{A(u)+H_\theta}+\mu_{\mathrm{lat}})
                                             \,\dd\gamma_n(u).
                                  \label{eq:external-matter-factor}
\end{align}
Then \eqref{eq:external-zero-sector} and
\eqref{eq:external-Gaussian-change} imply
\begin{equation}\label{eq:external-zero-factorization}
 Z_n^{(0)}[1]=Z_n^{\mathrm{ph}}z_{r,n},\qquad
                         0<z_{r,n}\le K_n(0).
\end{equation}
Positivity follows from that of $Z_n^{(0)}[1]$;
the upper bound follows because both integration
measures in \eqref{eq:external-matter-factor} are
probability measures.

On a cubic torus $V=n^3$, fix a positive reference mass
$\bar m$, mesh $\eta$ and a real pair
$(m_{\mathrm{in}},\varepsilon_{\mathrm{in}}^{\mathrm v})$,
with $\mu_{\mathrm{lat}}=\bar m\eta+m_{\mathrm{in}}$.
The density with external cochain $a$ is
\begin{equation}\label{eq:external-matter-density}
 \mathrm{dens}^{\theta,a}(u)=
 \mathrm e^{-\varepsilon_{\mathrm{in}}^{\mathrm v}n^3}
 \frac{\det_n(D_{A(u)+H_\theta+a}+\mu_{\mathrm{lat}})}
                       {\det_n(D^{\mathrm{ch}}[0]+\bar m\eta)}.
\end{equation}
The denominator is positive, since the real part
of $D^{\mathrm{ch}}[0]+\bar m\eta$ is at least $\bar m\eta>0$.
At $a=0$ this is a nonnegative bounded continuous
density. Let $\nu^{\theta,0}$ be its measure against
$\gamma_n$ and let $\bar\nu$ be the twist average.
Equations \eqref{eq:external-zero-sector} and
\eqref{eq:external-Gaussian-change} give
\begin{align}
 \bar\nu(\Phi)
   &=\frac{\mathrm e^{-\varepsilon_{\mathrm{in}}^{\mathrm v}n^3}}
        {\det_n(D^{\mathrm{ch}}[0]+\bar m\eta)Z_n^{\mathrm{ph}}}
                                      Z_n^{(0)}[\Phi],
                                  \label{eq:external-density-zero-sector}\\
 \bar\nu(1)
   &=\frac{\mathrm e^{-\varepsilon_{\mathrm{in}}^{\mathrm v}n^3}
                                   z_{r,n}}
                       {\det_n(D^{\mathrm{ch}}[0]+\bar m\eta)}
                                  \in(0,\infty).
                                  \label{eq:external-density-normalization}
\end{align}
Thus $\bar\nu/\bar\nu(1)$ is exactly the zero-flux
curvature law. The vacuum normalization cancels in
this quotient. The quantities $Z_n^{\mathrm{ph}}$
and $z_{r,n}$ depend only on $(n,q,\mu_{\mathrm{lat}})$
and the fixed chamber parameters; the auxiliary
split $\mu_{\mathrm{lat}}=\bar m\eta+m_{\mathrm{in}}$
does not change them.

\subsection{Holomorphic external partition functions}
\label{subsec:external-holomorphy}

Write $\mathfrak F(n)$ for the exponentially bounded
Borel functions $\Phi$ used above. For $w\in\Z^3$ let
$\mathcal M_w$ be the closed integer two-cochains of
flux $w$. For a complex one-cochain $a$ define
\begin{align}
 Z_n^{\mathrm{x},(w)}(a;\Phi)
  &=\sum_{z\in\mathcal M_w}\int_{[0,r_q)^{\mathcal L_n}}
       \Phi(F)\mathrm e^{-\|F\|^2/2}
                \det_n(D_{A+a}+\mu_{\mathrm{lat}})\,\dd A,
       \qquad F=dA+r_qz,\label{eq:external-sector-source}\\
 \mathbf Z_n^{\mathrm{x}}(a;\Phi)
  &=\int\frac{\dd\theta}{(2\pi)^3}
       \int_{V_{\mathrm t}}\Phi(dA)\mathrm e^{-\|dA\|^2/2}
       \det_n(D_{A+H_\theta+a}+\mu_{\mathrm{lat}})\,\dd\sigma(A).
                                  \label{eq:external-averaged-source}
\end{align}
When $\Phi=1$ we omit it from the notation. These are
partition functions in an external connection; the
curvature $F$ in their photon weight is not shifted.

\begin{proposition}\label{prop:external-entire}
For every $\Phi\in\mathfrak F(n)$ both expressions
converge absolutely, locally uniformly in $a$, and are
entire in its $3V$ complex link values. The sum over
all sectors converges absolutely, locally uniformly
as well. Moreover
\begin{equation}\label{eq:external-source-zero-sector}
 Z_n^{\mathrm{x},(0)}(a;\Phi)
             =\mathcal N_n\mathbf Z_n^{\mathrm{x}}(a;\Phi),
 \qquad
 Z_n^{\mathrm{x},(w)}(0;\Phi)=Z_n^{(w)}[\Phi].
\end{equation}
In particular $\mathbf Z_n^{\mathrm{x}}(0)>0$.
\end{proposition}

\begin{proof}
For real $A,\theta$ the imaginary parts of $A+a$ and
$A+H_\theta+a$ are both $\operatorname{Im}a$. On a
polydisc $\max_b|a(b)|\le s$ the determinant is
bounded by $K_n(s)$. If
$|\Phi(u)|\le C_\Phi \mathrm e^{C_\Phi\|u\|}$, the absolute
value of each integrand is consequently bounded by
\[
 K_n(s)C_\Phi \mathrm e^{C_\Phi^2}\mathrm e^{-\|F\|^2/4}
 \quad\text{or}\quad
 K_n(s)C_\Phi \mathrm e^{C_\Phi^2}\mathrm e^{-\|dA\|^2/4}.
\]
The first majorant can be summed over every integer
two-cochain and integrated over the compact box,
giving
$K_n(s)C_\Phi \mathrm e^{C_\Phi^2}(r_qC_q(1/4))^{3V}$.
The second is integrable by
\eqref{eq:external-transverse-Gaussian}: rescaling
$A$ by $\sqrt2$ gives
\[
 \int_{V_{\mathrm t}}\mathrm e^{-\|dA\|^2/4}\,\dd\sigma(A)
                          =2^{V-1}\mathcal G_n.
\]
These bounds also show local uniform absolute
convergence of the sector sum.

We recall the parameter-integral argument to justify
entireness and all later differentiations. Let
$(Y,\nu)$ be a $\sigma$-finite measure space. Suppose
$v(y,a)$ is measurable in $(y,a)$, entire in
$a\in\C^d$ for each $y$, and bounded on each
polydisc by a $\nu$-integrable $M_s(y)$. If $v_\alpha(y)$
are its Taylor coefficients at zero, integration
over the product of circles of radius $s$ gives
\[
 v_\alpha(y)=s^{-|\alpha|}
       \int_{[0,2\pi]^d}
       v(y,s\mathrm e^{i\theta})\mathrm e^{-i\alpha\cdot\theta}
                         \frac{\dd\theta}{(2\pi)^d},
 \qquad |v_\alpha(y)|\le s^{-|\alpha|}M_s(y).
\]
The integral formula makes the coefficients measurable.
For $\max_j|a_j|\le s$, applying the bound at radius
$2s$ yields
\[
 \sum_\alpha\int|v_\alpha(y)a^\alpha|\,\dd\nu(y)
       \le\sum_\alpha2^{-|\alpha|}\int M_{2s}(y)\,\dd\nu(y)
       =2^d\int M_{2s}(y)\,\dd\nu(y)<\infty.
\]
Thus the Taylor series may be integrated term by term,
and its integral is an entire function with coefficients
$\int v_\alpha$. This also proves differentiation under
the integral at zero. Translating the parameter to an
arbitrary $a^0$ proves it there, using the majorant
$M_{s+\|a^0\|_\infty}$. This argument applies to a
counting measure times Lebesgue measure as well as to
the transverse integral.

Apply this argument to the two displayed majorants.
The dependence on $a$ in each integrand is a finite
sum of exponentials of linear forms, by the determinant
polynomial. The other factors are Borel and independent
of $a$, so the hypotheses hold also when $\Phi$ is
merely Borel. This proves the asserted holomorphy.
For fixed complex $a$, the function
$W_a(A)=\det_n(D_{A+a}+\mu_{\mathrm{lat}})$ on real
cochains is bounded by $K_n(\|\operatorname{Im}a\|_\infty)$,
continuous, and has the two invariances required in
\eqref{eq:external-zero-unfolding}. That formula gives
\eqref{eq:external-source-zero-sector}; the value at zero
follows from the definitions. Its positive zero value
follows from $Z_n^{(0)}[1]>0$ and $\mathcal N_n>0$.
\end{proof}

The complex densities in
\eqref{eq:external-matter-density} define finite complex
measures $\nu^{\theta,a}$ and their twist average
$\bar\nu^a$ against $\gamma_n$. Their absolute values
are bounded by
\begin{equation}\label{eq:external-complex-density-bound}
 |\mathrm{dens}^{\theta,a}(u)|
 \le
 \frac{\mathrm e^{-\varepsilon_{\mathrm{in}}^{\mathrm v}n^3}
                   K_n(\|\operatorname{Im}a\|_\infty)}
                   {\det_n(D^{\mathrm{ch}}[0]+\bar m\eta)}.
\end{equation}
This bound is independent of $(\theta,u)$.
Fubini therefore writes the averaged measure as a
bounded Borel density against $\gamma_n$. Applying
\eqref{eq:external-Gaussian-change} at each twist,
and then \eqref{eq:external-source-zero-sector}, gives
\begin{equation}\label{eq:external-complex-measure}
 \bar\nu^a(\Phi)=
 \frac{\mathrm e^{-\varepsilon_{\mathrm{in}}^{\mathrm v}n^3}}
      {\det_n(D^{\mathrm{ch}}[0]+\bar m\eta)\mathcal G_n}
                     \mathbf Z_n^{\mathrm{x}}(a;\Phi)
 =
 \frac{\mathrm e^{-\varepsilon_{\mathrm{in}}^{\mathrm v}n^3}}
      {\det_n(D^{\mathrm{ch}}[0]+\bar m\eta)Z_n^{\mathrm{ph}}}
                     Z_n^{\mathrm{x},(0)}(a;\Phi).
\end{equation}
The exponential majorants in the proof above justify
both changes of integration for $\Phi\in\mathfrak F(n)$.
In particular, for every Borel set $E\subset F_n$,
$a\mapsto\bar\nu^a(E)$ is entire and
\[
 |\bar\nu^a(E)|\le
 \frac{\mathrm e^{-\varepsilon_{\mathrm{in}}^{\mathrm v}n^3}
                   K_n(\|\operatorname{Im}a\|_\infty)}
                   {\det_n(D^{\mathrm{ch}}[0]+\bar m\eta)}.
\]
Indeed $F_n$ is a closed linear subspace, so its
Borel set indicator extends to a bounded Borel
function on $C^2(n)$. Composition with any linear
parameter map $t\mapsto a(t)$ remains entire.
These statements concern the unnormalized complex
measures; the positive normalization at $a=0$ is
\eqref{eq:external-density-normalization}.

\subsection{The source shift and Wick derivatives}
\label{subsec:external-source-shift}

\begin{proposition}\label{prop:external-shift}
For every flux $w$ and complex one-cochain $a$,
\begin{equation}\label{eq:external-shift-sector}
 Z_n^{(w)}[\mathrm e^{\langle da,F\rangle}]
      =\mathrm e^{\langle da,da\rangle/2}Z_n^{\mathrm{x},(w)}(a).
\end{equation}
Both sides, and their sums over $w$, converge
absolutely. Consequently
\begin{equation}\label{eq:external-shift-all-flux}
 Z_n[\mathrm e^{\langle da,F\rangle}]
      =\mathrm e^{\langle da,da\rangle/2}
                    \sum_w Z_n^{\mathrm{x},(w)}(a).
\end{equation}
The quadratic form in the prefactor is bilinear;
there is no complex conjugation of $da$.
\end{proposition}

\begin{proof}
First let $a$ be real. On real cochains define
\[
 P_a(A)=\sum_{z\in\mathcal M_w}
  \mathrm e^{-\|dA+r_qz\|^2/2+\langle da,dA+r_qz\rangle}
                            \det_n(D_A+\mu_{\mathrm{lat}}).
\]
This is a nonnegative Borel function. Adding
$r_q\mathbf1_b$ to $A$ adds $d\mathbf1_b$ to the
integer two-cochain in its curvature. This map
permutes $\mathcal M_w$: its image is closed
and has the same flux, since each exact plaquette
cochain has zero plane sums. The determinant
is unchanged. Hence $P_a$ is invariant under
$r_qC^1(n;\Z)$. Both the compact box and its
translate by $a$ are fundamental domains for
this group. The domain argument in
\cref{subsec:external-zero-flux} gives
\[
 \int_{[0,r_q)^{\mathcal L_n}}P_a(A)\,\dd A
        =\int_{[0,r_q)^{\mathcal L_n}}P_a(A+a)\,\dd A.
\]
On the right the exponent completes to
\[
 -\tfrac12\|F+da\|^2+\langle da,F+da\rangle
                  =-\tfrac12\|F\|^2+\tfrac12\|da\|^2.
\]
Tonelli now identifies the two sides with
\eqref{eq:external-shift-sector} for real $a$.
Their finiteness also follows from the Gaussian
majorant used below.

For complex $a$, the right side of
\eqref{eq:external-shift-sector} is entire by
\Cref{prop:external-entire}.
The left side is entire by the same
parameter-integral argument. In fact, on
$\max_b|a(b)|\le s$,
\[
 \|d\operatorname{Re}a\|\le4s\sqrt{3V}=:c_s,
 \qquad
 \left|\mathrm e^{-\|F\|^2/2+\langle da,F\rangle}
                            \det_n(D_A+\mu_{\mathrm{lat}})\right|
                  \le K_n(0)\mathrm e^{c_s^2}\mathrm e^{-\|F\|^2/4}.
\]
This can be summed and integrated over all integer
two-cochains, so it supplies an integrable polydisc
majorant. Both sides agree on the real cochain
space. An entire function vanishing there has
all complex partial derivatives zero there:
at a real point these derivatives can be taken
along real increments. Repeating this for all
derivatives shows that every Taylor coefficient
at zero vanishes. Thus the two entire functions
agree for all complex $a$.
Finally the same Gaussian majorant permits
summation over sectors; the sector sum on the
right converges absolutely by
\Cref{prop:external-entire}.
\end{proof}

Let $a^{(1)},\ldots,a^{(r)}$ be real cochains and
$X_j=\langle da^{(j)},F\rangle$. Dividing
\eqref{eq:external-shift-sector} at $w=0$ by its
positive zero-source value yields the normalized
generating identity
\begin{equation}\label{eq:external-generating-general}
 \E_n^{\mathrm{av}}\exp\left(
       \sum_jz_jX_j-\tfrac12\sum_{i,j}
                  z_iz_j\langle da^{(i)},da^{(j)}\rangle\right)
 =
 \frac{\mathbf Z_n^{\mathrm{x}}(\sum_jz_ja^{(j)})}
                         {\mathbf Z_n^{\mathrm{x}}(0)},
                                                   \qquad z\in\C^r.
\end{equation}
The expectation $\E_n^{\mathrm{av}}$ is the zero-flux
law identified in \eqref{eq:external-zero-sector}.
The Gaussian factor on the left is determined by
the counting pairing of the curvatures of the
source cochains.

Use the reference cells of \Cref{sec:reference-cells}
on the torus $n_*$. Their index set is
$\mathcal C_*=\{0,\ldots,2L-1\}^3$, and
$\widehat a_c$ is the periodization of
$\hat{\mathfrak a}_c^o$. These are real periodic cochains.
Equation \eqref{eq:cells-orthogonality} gives
\begin{equation}\label{eq:external-cell-pairing}
 \langle d\widehat a_c,d\widehat a_{c'}\rangle
                         =v_\eta\,1_{\{c=c'\}}.
\end{equation}
The variance $v_\eta=\|d\hat{\mathfrak a}_\eta^o\|^2$
is independent of the cell and the placement.
This is the cell-geometry input needed here:
the external-response proof uses neither the
cell width nor its offsets otherwise. Put
\[
 a(t)=e^{-1}\sum_{c\in\mathcal C_*}t_c\widehat a_c,\qquad
 \bar Z^{\mathrm{ext}}(t)=
     \int\frac{\dd\theta}{(2\pi)^3}
        \int\det_{n_*}(D_{A(u)+H_\theta+a(t)}
                              +\mu_{\mathrm{lat}})\,\dd\gamma_{n_*}(u).
\]
The reference torus has $n_*\ge4$, and $e>0$ is
the physical charge in this convention.
Equation \eqref{eq:external-Gaussian-change} gives
\begin{equation}\label{eq:external-reference-partition}
 \bar Z^{\mathrm{ext}}(t)=
          \mathcal G_{n_*}^{-1}\mathbf Z_{n_*}^{\mathrm{x}}(a(t)),
 \qquad
 \mathring\Phi(t):=\frac{\bar Z^{\mathrm{ext}}(t)}
                           {\bar Z^{\mathrm{ext}}(0)}
        =\frac{\mathbf Z_{n_*}^{\mathrm{x}}(a(t))}
                           {\mathbf Z_{n_*}^{\mathrm{x}}(0)}.
\end{equation}
Its zero value in the denominator is positive,
and $\mathring\Phi$ is entire. With
$X_c=\langle d\widehat a_c,F\rangle$,
\eqref{eq:external-generating-general} becomes
\begin{equation}\label{eq:external-generating-cells}
 \E_{n_*}^{\mathrm{av}}
      \exp\sum_c(z_cX_c-\tfrac12v_\eta z_c^2)
                                      =\mathring\Phi(ez).
\end{equation}

Define the Hermite polynomials with this variance by
$H_j(y)=\partial_z^j \mathrm e^{zy-v_\eta z^2/2}|_{z=0}$.
In particular
\[
 H_2(y)=y^2-v_\eta,\qquad
 H_4(y)=y^4-6v_\eta y^2+3v_\eta^2.
\]
To differentiate \eqref{eq:external-generating-cells}
under the expectation, on a polydisc
$\max_c|z_c|\le s$ use
\[
 \left|\exp\sum_c(z_cX_c-\tfrac12v_\eta z_c^2)\right|
 \le
 \exp\left(s|\mathcal C_*|\sqrt{v_\eta}\|F\|
                       +\tfrac12|\mathcal C_*|s^2v_\eta\right).
\]
This is exponentially bounded in the real curvature
and is integrable under the zero-flux law, by the
Gaussian majorant. The parameter-integral argument
therefore applies to this probability measure.
The resulting products of $H_2$ or $H_4$ are also
integrable: $|X_c|\le\sqrt{v_\eta}\|F\|$ bounds them by a
finite polynomial in $\|F\|$, which is bounded by
a constant times $\mathrm e^{\|F\|}$.

For $m=2$ or $4$ and a subset $S\subset\mathcal C_*$,
differentiate in precisely those variables. The
factor $ez$ on the right supplies $e^{m|S|}$, giving
\[
 \E_{n_*}^{\mathrm{av}}\prod_{c\in S}e^{-m}H_m(X_c)
         =\left(\prod_{c\in S}\partial_{t_c}^m
                                      \mathring\Phi\right)(0).
\]
Expanding the product over subsets proves, for
every real $c_m$,
\begin{equation}\label{eq:external-Wick-product}
 \E_{n_*}^{\mathrm{av}}
        \prod_{c\in\mathcal C_*}(e^{-m}H_m(X_c)-c_m)
     =\left(\prod_{c\in\mathcal C_*}
                (\partial_{t_c}^m-c_m)\mathring\Phi\right)(0).
\end{equation}
The derivative operators commute because
$\mathring\Phi$ is entire. This identity is the
finite-volume quantity to which the terminal
interacting estimates will be applied; all its
normalization factors have already been included
in \eqref{eq:external-reference-partition}.


\section{Reference scales and rescaled probes}
\label{subsec:interacting-stop-probe}

We use the reference torus of the tuned flow and prove the
uniform probe estimates from \cite[\S21.1]{QED3Master}.
Fix $\bar m\in(0,1]$, put
$h=L^{\lfloor\log_L(1/\bar m)\rfloor}$, and assume
$0<e\sqrt h\le e_1$ and $N\ge N_{\rm ref}(e)$.
By \cite[(47)]{CoverTuning} and
\cite[(53)]{CoverTuning},
\begin{equation}\label{eq:interacting-physical-scales}
 \begin{aligned}
 \eta=L^{-N},\qquad K=N+\log_Lh,\qquad
 &a_j=L^{j-N},\\
 e_j=e\sqrt{a_j},\qquad
 &\eta_j=L^{-j}=\eta/a_j,\qquad \ell_j=R^*/a_j .
 \end{aligned}
\end{equation}
Thus a point with level-$j$ coordinates $c$
has physical coordinates $a_jc$, and
$T_j=\eta_j\Z^3/\ell_j\Z^3$.
At $j=0$ the charge is $q=e\sqrt\eta$.
The symbols $e_0,e_1$ without a running
index retain their meanings as charge ceilings.
Set $v_j=-\log e_j$, $p_{0,j}=v_j^{p_0}$ and $p_j=v_j^p$;
$r_j$ is the least nonnegative integer power of $L$ at least
$v_j^{r_{\rm loc}}$.
For a bond cochain $a$ on the fine unit torus,
its rescaling is
\begin{equation}\label{eq:interacting-connection-rescaling}
                    a^{[j]}(b)=L^{j/2}a(L^jb).
\end{equation}
The physical scale $a_j$ is a number, while
$a$ denotes a connection.
On the reference torus,
\cite[(54)]{CoverTuning} and
\cite[(50)]{CoverTuning} give
\[
 k_*\in\{K-1,K-2\},\qquad
 a_*\in\{h/L,h/L^2\},\qquad
 \ell_*=Mr_*,\qquad r_*>31L+16.
\]
The offset $K-k_*$ and the quantities $a_*,\ell_*,r_*$ are
independent of $N\ge N_{\rm ref}(e)$; the index $k_*$
depends on $N$ through $K$.
Write $r_{\rm loc}$ for the localization
power in \cite[\S6.1]{CoverTuning}.

Let $\beta\in C_c^\infty(\R^3;\Lambda^2)$ be real
and nonzero. Use the cube
$Q_\beta=z_\beta+[-D_\beta^{\rm buf}/2,D_\beta^{\rm buf}/2]^3$
of \Cref{sec:reference-cells},
where $D_\beta^{\rm buf}=D_\beta+4$.
In this section a placement $v\in\eta\Z^3$
is admitted if, in physical coordinates,
\begin{equation}\label{eq:interacting-admitted-placement}
              v+Q_\beta\subset[\ell_c/8,7\ell_c/8]^3,
              \qquad \ell_c=R^*/(2L).
\end{equation}
A centered placement has
$|v+z_\beta-(\ell_c/2)(1,1,1)|_\infty\le\eta/2$.
It is admitted when
$D_\beta^{\rm buf}/2+\eta/2\le3\ell_c/8$.
This follows by adding the displacement
of the cube center to its half-side.

Form $\widehat{\mathfrak a}$ with $\tau_v\beta$
as in \eqref{eq:cells-midpoint-source}, and
periodize it on $C(n_*)$, $n_*=R^*/\eta$.
For $c\in\mathcal C_*=\{0,\ldots,2L-1\}^3$,
write $\widehat{\mathfrak a}_c$ for its
cell copy; thus $\widehat{\mathfrak a}_c=\widehat a_c$ in the
introduction. In physical coordinates the
wall map carrying the first cell to cell $c$
acts in coordinate $\mu$ by translation through
$2i\ell_c$ if $c_\mu=2i$, and by reflection
$y_\mu\mapsto(2i+2)\ell_c-y_\mu$ if $c_\mu=2i+1$.
Define the marked connections by
\begin{equation}\label{eq:interacting-marked-connections}
 a(t)=e^{-1}\sum_ct_c\widehat{\mathfrak a}_c,\qquad
 \widehat c_c=e^{-1}\widehat{\mathfrak a}_c^{[k_*]},\qquad
 \mathscr C(t)=a(t)^{[k_*]}=\sum_ct_c\widehat c_c .
\end{equation}
For a level-$j$ bond function $A$, $\partial A$
uses forward differences divided by $\eta_j$,
and
\[
 |\delta_\alpha\partial A|_\infty
  =\sup_{\mu,\nu}\sup_{0<d(x,x')\le1}
     \frac{|\partial_\nu A_\mu(x')-\partial_\nu A_\mu(x)|}
          {d(x,x')^\alpha}.
\]
The external norm of a fine connection is
$|a|_{\rm ext}=|a^{[k_*]}|_\infty+
|\partial a^{[k_*]}|_\infty+
|\delta_\alpha\partial a^{[k_*]}|_\infty$.

We define the uniform probe constants over both
possible stop scales. For $\eta'=L^{-N'}$ put
\[
 g_\nu^{\eta'}[\beta']
      =-\sum_{\mu\ne\nu}\delta_{\eta',\mu}J_{\eta'}\beta'_{\mu\nu}.
\]
This is the smooth vector field
$J_{\eta'}a'_{\eta'}$ in the midpoint formula.
Let $\mathcal I_\beta$ consist of triples
$(N',a_{\rm s},r')$ with
$N'\ge0$, $a_{\rm s}\in\{h/L,h/L^2\}$,
$N'+\log_La_{\rm s}\ge1$, and
$r'\in\{\pm1\}^3$.
$r'$ is identified with its diagonal orthogonal
matrix, and $U_{r'}\beta(x)=(\wedge^2r')\beta(r'^{-1}x)$.
At such a triple set $\eta''=\eta'/a_{\rm s}$,
$g^0=g^{\eta'}[U_{r'}\beta]$, and
\[
 G_\nu(c)=\sqrt{a_{\rm s}}\,
                  g_\nu^0(a_{\rm s}c+\eta'e_\nu/2)
                       \quad(c\in\eta''\Z^3),\qquad
 W_\nu(x)=\sqrt{\eta'}\,g_\nu^0(\eta'x+\eta'e_\nu/2)
                       \quad(x\in\Z^3).
\]
The three norms of $G$ are taken at mesh $\eta''$.
Set
\begin{equation}\label{eq:interacting-uniform-probe-constants}
 C_\beta=\sup_{\mathcal I_\beta}
              (|G|_\infty+|\partial G|_\infty+
                                    |\delta_\alpha\partial G|_\infty),
 \qquad
 \mathbf D=\sup_{\mathcal I_\beta}\|dW\|_{\ell^2(\text{plaquettes})}^2.
\end{equation}
These suprema depend on $\bar m,L,\alpha,\beta$,
and contain no charge, torus, or placement.
They have the finite bounds
\begin{equation}\label{eq:interacting-uniform-probe-finite}
 C_\beta\le2\sqrt h(1+h+3h^2)|\beta|_{C^3},
 \qquad
 \mathbf D\le48|\beta|_{C^2}^2(D_\beta^{\rm buf}+3)^3 .
\end{equation}

Indeed the centered difference is the average
of the corresponding derivative along its
segment, and $J_{\eta'}$ has norm one on
each derivative sup norm. Each component
of $g^0$ has two summands, hence
$|D^ig_\nu^0|_\infty\le2|\beta|_{C^{i+1}}$ for
$i=0,1,2$. This gives
$|G|_\infty\le2a_{\rm s}^{1/2}|\beta|_{C^1}$.
A first lattice difference of $G$ is
$a_{\rm s}^{3/2}$ times an average of
$\partial_\mu g_\nu^0$, so
$|\partial G|_\infty\le2a_{\rm s}^{3/2}|\beta|_{C^2}$.
The difference of these averages at $c,c'$
is bounded by
$6a_{\rm s}|c-c'|\,|\beta|_{C^3}$.
Since $|c-c'|\le1$ and $\alpha<1$,
the third norm is at most
$6a_{\rm s}^{5/2}|\beta|_{C^3}$.
Sum these bounds and use $a_{\rm s}\le h$.

For the energy, each curl component of $W$
is a difference of two first differences,
and has modulus at most
$4(\eta')^{3/2}|\beta|_{C^2}$.
A supporting plaquette has a supporting
bond, whose endpoints in physical coordinates
lie in $r'Q_\beta$.
Its basepoint is therefore in that cube
enlarged by $\eta'$ in each negative
coordinate direction. There are at most
$(D_\beta^{\rm buf}/\eta'+2)^3$ basepoints
and three plaquette directions. Squaring
the pointwise curl bound gives
$48|\beta|_{C^2}^2(D_\beta^{\rm buf}+2\eta')^3$,
which implies \eqref{eq:interacting-uniform-probe-finite}.

At the reference stop, $(N,a_*,1)\in\mathcal I_\beta$.
The midpoint formula and
\eqref{eq:interacting-connection-rescaling} give
\[
 \widehat{\mathfrak a}^{[k_*]}_\nu(c)
       =\sum_{z\in\Z^3}G_\nu(c-v/a_*+\ell_*z).
\]
Supporting bond endpoints lie in
$a_*^{-1}(v+Q_\beta)$.
Distinct periodic copies of that cube have
distance at least
$\ell_*-3\ell_*/(8L)>\ell_*/2>1+2\eta_*$.
Each functional defining one of the three
norms reads only bonds with endpoint distances
at most $1+2\eta_*$. It therefore reads
at most one periodic copy. Taking its
supremum identifies each torus norm with
the corresponding norm of $G$.
Translations and coordinate reflections
preserve these norms, including the one-bond
shift that reverses an oriented component.
The same argument applies to the cell copies.

The image of $v+Q_\beta$ under each wall map
lies in
$\prod_\mu[c_\mu\ell_c+\ell_c/8,(c_\mu+1)\ell_c-\ell_c/8]$.
For distinct cells, a separating coordinate
on the periodic circle gives distance
at least $\ell_c/4$: any joining arc
contains the two margins of size $\ell_c/8$.
Thus supporting endpoints of distinct
$\widehat c_c$ have stop-coordinate distance
at least $\ell_*/(8L)$, which is larger
than $1+2\eta_*$. Again each external-norm
functional reads at most one copy, so
\begin{equation}\label{eq:interacting-external-probe-bound}
 |a(t)|_{\rm ext}\le C_\beta|t|_\infty/e,\qquad
                 v_\eta=\|d\widehat{\mathfrak a}\|^2\le\mathbf D.
\end{equation}
For the second bound, the covering curvature
$dW$ is supported in the cube of side
$D_\beta^{\rm buf}/\eta+2<n_*$
obtained by enlarging $\eta^{-1}(v+Q_\beta)$.
Its periodic translates do not overlap,
so periodization preserves its counting
norm. Distinct cell curvatures also have
disjoint plaquette supports: a common
plaquette would have vertices in both
endpoint sets at physical distance at
most $\eta\le1$, whereas $\ell_c/4>1$.
This gives \eqref{eq:cells-orthogonality}
for every admitted placement.

Let $P_c$ be the unit-block centers whose
blocks contain an endpoint of a supporting
bond of $\widehat c_c$, and let $P_c^+$
be their $(M+1)$-neighborhood.
For $y\in P_c$, $y'\in P_{c'}$ choose the
corresponding endpoints in $\Delta_y,\Delta_{y'}$.
Each endpoint is at distance less than $1/2$
from its center. Consequently
\begin{align}
 d_{\rm sep}:=\min_{c\ne c'}d(P_c,P_{c'})
    &\ge\ell_*/(8L)-1=Mr_*/(8L)-1,
                              \label{eq:interacting-probe-separation}\\
 |P_c^+|&\le p_\beta:=
                  (D_\beta^{\rm buf}/a_*+2M+4)^3 .
                              \label{eq:interacting-support-number}
\end{align}
For the count, the endpoints have a lift in
a cube of side $D_\beta^{\rm buf}/a_*$.
A center adds less than $1/2$ to each
coordinate bound, and its neighborhood adds
$M+1$. The resulting open interval has
length $D_\beta^{\rm buf}/a_*+2M+3$
and contains fewer than one plus that
many integers. Multiply the three counts.
If the lattice probe vanishes, all these
sets are empty and $d_{\rm sep}=\infty$.

\section{The terminal inputs}\label{sec:rrm-terminal-inputs}

We recall the terminal estimates used in the response proof, with the
normalizations of \cite[\S6 and Appendix~A]{TerminalPressure}.
They apply to the common trajectory fixed in the introduction.
The marked activity estimates, logarithm and complement argument
are proved in the following sections.

\subsection{Operators, norms and constants}
At the reference stop $T_* =\eta_*\Z^3/\ell_*\Z^3$,
coarse spinors live at $T_*^0=\Z^3/\ell_*\Z^3$ with weight one;
fine chamber indices have weight $\eta_*^3/48$.
The block average $Q_*$ and cluster injection $J$ are those of
\cite[\S2.1 and \S4.2]{FermionBlock}. The fine resolvent and blocked operator are
\begin{equation}\label{eq:interacting-terminal-operators}
 \begin{split}
 F_*(A)&=D_A+\mu_*+b_*JP_*(A)J^*,\qquad S_*(A)=F_*(A)^{-1},\\
 H_*(A)&=b_*S_*(A)JQ_*^T(-A),\qquad
 D_*(A)=b_*-b_*Q_*(A)J^*H_*(A).
 \end{split}
\end{equation}
Here $P_*=Q_*^T(-A)Q_*(A)$ on the cluster-averaged space,
and $H_*^\#$ includes ordinary and covariant H\"older-pair rows.
The barred rows are their charge-conjugate partners.
The parameter ranges are
\[
 b_-=1-L^{-1}\le b_*\le1,\qquad
 \mu_-=L^{-3}<\mu_*\le L^{-1},\qquad
 \gamma^FM\ge16,\qquad \ell_*>2M+2.
\]
The positive mass $\mu_*=\bar m a_*$ is distinct from the running
mass shift $m_*$ and from the possibly signed total fine mass.

The domains $\widetilde R_j$, $\mathfrak X_j$ and
$\widetilde G_j$ are the background, external and total-connection
domains of \cite[\S6]{MaxwellBlock} and \cite[\S2.3]{NormalizedContraction}.
In particular
\[
 \mathfrak X_j=\{C:|C|_\infty<e_j^{-1+2\varepsilon},\
                   |\partial C|_\infty<e_j^{-1+3\varepsilon}\},
 \qquad \widetilde R_j+\mathfrak X_j\subset\widetilde G_j.
\]
The local real-gauge definition of $\widetilde R_j$ makes it invariant
under real constants. At the stop put
\[
 \mathcal W=\{A_r+C:A_r\in\widetilde R_*\text{ real},\
                  C\in\mathfrak X_*\}.
\]
On this domain the inverse and field maps are analytic. Their
localized majorants, including short-pair H\"older rows, are bounded
by $C^F E(y,y')$, where $E(y,y')=\mathrm e^{-\gamma^Fd(y,y')}$.
For long pairs the distance is the minimum of the distances from
the two endpoint blocks. All global scalar component row sums of
$H_*^\#$ and its barred partner are at most $C^F$.

Localize a fine index in its unit block and a pair at its first
point. For an operator $O$ define
\begin{align*}
 |O|(y,y')=\sup_{\|v\|_\infty\le1}
                       \|\mathbf1_yO\mathbf1_{y'}v\|_\infty,\qquad&\\
 \|O\|_\gamma=\max\left\{
 \sup_y\sum_{y'}|O|(y,y')\mathrm e^{\gamma d(y,y')},\
 \sup_{y'}\sum_y|O|(y,y')\mathrm e^{\gamma d(y,y')}\right\}.&
\end{align*}
The spinor norm here is the maximum component modulus.
For coarse operators, $\|O\|_{\mathrm{op}}\le2\|O\|_0$ and
$\|O\|_{\mathfrak S_1}\le4\sum_{y,y'}|O|(y,y')$, where
$\mathrm{op}$ and $\mathfrak S_1$ denote the Hilbert-space operator
norm and trace norm, respectively.

For a direction $B$, let $P_B$ be the unit blocks containing
an endpoint of a bond in its support, put
$P_B^+=\{y:d(y,P_B)\le M+1\}$ and
$\epsilon_B=e_*|B|_\infty$.
The first coarse derivative has majorant
$226(C^F)^2\epsilon_B\sum_{z\in P_B}E(y,z)E(z,y')$.
The first field-leg derivative has the analogous two-edge chain
through $P_B^+$, with coefficient $C_H\epsilon_B$ and a sum
over at most two endpoint blocks. For directions separated by
$M+2$, both mixed derivatives have the sum of the two ordered
three-edge chains through $P_B,P_{B'}$, or their enlarged sets
for the legs, with coefficient $C_H^2\epsilon_B\epsilon_{B'}$.
These precise pointwise estimates are
\cite[Proposition~A.1]{TerminalPressure}; the response proof uses
their support factors as well as their row sums.

We retain their constants:
\begin{align*}
 \mathcal S&=C^F(24/\gamma^F)^3,&
 h_{\rm L}&=6\mathcal S,& C_H&=203M(C^F)^2,\\
 C_K&=226\mathcal S^2,& C_K'&=1328\mathcal S^2,&
 c_\tau&=\min\{1/8,\mu_-/48,(292\mathcal S)^{-1}\},\\
 g&=(b_-^{-1}+\mu_-^{-1})^{1/2},&h_G&=1+g,&g_1&=b_-^{-1}+4/\mu_-.
\end{align*}
Set $\theta_{\rm CT}=\mu_-/(12\mathrm e)$ and
\[
 \vartheta_0=\min\{\theta_{\rm CT}/2,\gamma^F/8,1/M\},\qquad
 \alpha_G=2b_-^{-1}
           +4\mathrm e\mu_-^{-1}(24/\theta_{\rm CT})^3.
\]
The real covariance $G=D_*(A_r)^{-1}$, for every real $A_r$,
satisfies $\|G\|_{\mathrm{op}}\le g^2$ and has weighted component row and
column sums at most $\alpha_G$, with weight
$\mathrm e^{\vartheta_0d(u,u')}$.
On $\mathcal W^+=\{A+A_1:A\in\mathcal W,\ e_*|A_1|_\infty\le c_\tau\}$,
the first derivative bound holds with $C_K'$.
If $A=A_r+A_1$, $A_r$ is real, $e_*|A_1|_\infty\le\mu_*/24$
and $F_*(A)$ is invertible, then $D_*(A)$ is invertible and
$\|D_*(A)^{-1}\|_{\mathrm{op}}\le g_1$.
These are \cite[(72)--(75)]{TerminalPressure}.

The activity space, independent ordinary and pair generators, and
their coefficient norm are those of \cite[\S2.2]{NormalizedContraction}.
An activity supported on an $M$-polymer $X$ has coefficients
$E_{nm}$, integrated with the fine-index weights above and pair
weights $(\eta_*^3/48)^2$, and norm
$\sum_{n,m}h_1^nh_2^m\|E_{nm}\|/(n!m!)$.
The family norm adds the factor $\mathrm e^{\kappa d_M(X)}$ and takes
the supremum over polymers and the declared connection domain.
Here $Md_M(X)$ is the infimum of the lengths of connected finite
nets meeting every closed $M$-cube of $X$. A net is a finite union
of images of closed straight segments in the continuum torus;
its length is the sum of their Euclidean lengths, with the infimum
also taken over segment representations. A point is allowed as a
segment. We use the construction constants
$\kappa_0=54\log72$, $K_{\mathrm{pol}}=72^{54}/36$,
$c_{\rm rb}=54K_{\mathrm{pol}}$ and $\kappa\ge7\kappa_0+7$.
The anchored polymer sum is bounded by $K_{\mathrm{pol}}$, and its reblocking
bound has constant $c_{\rm rb}$, as proved in
\cite[Appendix~B, (81) and (83)]{NormalizedContraction}.
The fixed constant $C_0\ge C^F$ bounds the photon kernel rows and
the field substitution; in particular at
$a_2=C_0^{-1}e_*^{-1/4+\varepsilon}$,
\[
 |E(H_*^\#\Psi)|_{a_2}\le\|E\|_{h_*},\qquad
 |\langle\bar\psi,\psi\rangle|_{a_2}
          \le4e_*^{-1/2+2\varepsilon}\ell_*^3,
 \quad h_*=(e_*^{-1/4},e_*^{-1/4+\varepsilon}).
\]
The scalar constant $C_d$ is the local determinant-activity bound
in \eqref{eq:interacting-determinant-input}.
All constants in this subsection are fixed before the physical charge
and are independent of the fine cutoff, twist and placement.

\subsection{The measure identity and Gaussian domination}
We use the exact terminal density of
\cite[Theorem~B]{TerminalPressure}, including its external slot.
Its original measure is $\nu^{\theta,a}$ from
\eqref{eq:external-matter-density}; write $\bar\nu$ for its
zero-source twist average and $\bar\mu=\bar\nu/\bar\nu(1)$.
The full photon-coordinate map sends an exact fine curvature
$F\in F_N=dC^1(C(n_*);\R)$ to $(A_*;\zeta_j)_{j<k_*}$.
Its Gaussian marginal at the stop is denoted $\mathsf m_*$.
Every scalar fluctuation coordinate is an orthonormal linear
functional of $F$; every coarse plaquette has row norm at most one.
The numbers of coordinates are at most
$2\ell_*^3L^{3(k_*-j)}$ at level $j$ and $3\ell_*^3$ at the stop.
The stop Gaussian is nondegenerate on the mean-zero coclosed space.

Set $\chi_*(A)=\mathbf1_{\{|dA|_\infty\le p_*\}}$ and
\[
 \mathfrak G=\{|\zeta_j(\Upsilon)|\le p_{0,j}\ (j<k_*),\
                                  |dA_*|_\infty\le p_*\}.
\]
For every bounded Borel function $\Phi$ of $A_*$,
\begin{equation}\label{eq:rrm-terminal-Borel}
 \nu^{\theta,a}(\mathbf1_{\mathfrak G}\Phi(A_*))
 =\int\Phi(A)\chi_*(A)
             Y_*^x(A+c_\theta;\mathscr C)\,\dd\mathsf m_*(A),
 \qquad \mathscr C=a^{[k_*]},\quad
 c_{\theta,\mu}=\theta_\mu/(e_*\ell_*).
\end{equation}
The density is jointly continuous on the indicated support and is
invariant under admissible real gauge changes. Its explicit fermion
formula is recalled in \cref{subsec:interacting-determinant-ratio}.
The source-free normalized measure satisfies
\begin{equation}\label{eq:rrm-zero-Gaussian}
 \bar\mu(\mathrm e^{\langle f,F\rangle})\le \mathrm e^{\|f\|^2/2}
                 \qquad(f\text{ a real two-cochain}).
\end{equation}
This is the zero-sector result
\cite[Proposition~2.1]{TerminalPressure}.
The tower and Borel identity are
\cite[Proposition~3.1, Theorem~B]{TerminalPressure}; they precede
the response argument and do not require its conclusion.
\section{The terminal determinant ratio}
\label{sec:interacting-response}

The argument in this and the following sections is the terminal
response proof of \cite[\S23]{QED3Master}.
The free coefficients do not by themselves control the
reference Wick products. We must compare the normalized
interacting generating function with its one-cell free
factors, including mixed derivatives in different cells.
The quantity to be estimated is
\begin{equation}\label{eq:interacting-reference-error}
 \Theta_m^{(0)}=
 \left|\left(\prod_{c\in\mathcal C_*}
       (\partial_{t_c}^m-c_m)\mathring\Phi\right)(0)\right|^{1/N_*},
 \qquad m=2,4,\qquad N_*=|\mathcal C_*|=8L^3,
\end{equation}
where $c_m=N_0^{(m)}(0)$. The exact relation to Wick
products is \eqref{eq:external-Wick-product}.
The proof uses decay between separate insertions,
an analytic logarithm for the terminal fermion integral,
and a separate estimate of the complement of the
small-field set.

\subsection{Analyticity and normalization}
\label{subsec:interacting-determinant-ratio}

For the marked connections of
\eqref{eq:interacting-marked-connections}, set
\begin{equation}\label{eq:interacting-mark-radii}
 C_b=1+\sqrt h\,C_\beta,\qquad
 t_x=\frac{e_*^{3\varepsilon}}{C_b},\qquad
 t_x^+=\frac{e_*^{3\varepsilon}}{\sqrt h\,C_\beta},
 \qquad r=t_x/4.
\end{equation}
If $C_\beta=0$, interpret $t_x^+=\infty$.
For $t$ in the polydisc
$\mathbb D(t_x^+)=\{|t_c|<t_x^+\text{ for every }c\}$,
all the rescaled connections $a(t)^{[j]}$
belong to $\mathfrak X_j$ and
$e_*|\mathscr C(t)|_\infty<e_*^{3\varepsilon}$.
Indeed, \eqref{eq:interacting-external-probe-bound}
and $e_*\le e\sqrt h$ give
$|a(t)|_{\rm ext}<e_*^{-1+3\varepsilon}$.
This bounds the stop sup and first-difference
norms by the defining thresholds of $\mathfrak X_*$.
For $i=k_*-j$ the rescaling multiplies
the sup norm by $L^{-i/2}$ and the first
differences by $L^{-3i/2}$, while
$e_j=L^{-i/2}e_*$. Since
$\varepsilon<1/3$, the rescaled bounds are
strictly below $e_j^{-1+2\varepsilon}$ and
$e_j^{-1+3\varepsilon}$, respectively.

We first control the determinant activities accumulated
along the flow. Recall the local determinant
function $E_j^d(X,A)$ of
\cite[Proposition~6.1]{FermionBlock} and set
$E_j^{\rm det}(X,A)=E_j^d(X,A)-E_j^d(X,0)$.
Its analyticity, locality, real gauge invariance,
and evenness will all be used below; the relevant
bound is
\begin{equation}\label{eq:interacting-determinant-input}
 |E_j^d(X,A)|\le C_dM^3\mathrm e^{-(\kappa+1)d_M(X)}
                       \qquad(A\in\widetilde G_j).
\end{equation}
For a level-$(j+1)$ bond cochain $C$, write
$C_L(b)=L^{-1/2}C(b/L)$ for its one-step rescaling.
Let $\mathfrak P_0=0$ and
\begin{equation}\label{eq:interacting-P-recursion}
 \mathfrak P_{j+1}(Y,C)=
 \sum_{\overline X^{\,L}=LY}
       \bigl(\mathfrak P_j(X,C_L)+E_j^{\rm det}(X,C_L)\bigr).
\end{equation}
Then $\mathfrak P_j(X,\cdot)$ is analytic on
$\mathfrak X_j$, vanishes at zero, and depends only
on bonds meeting $X$. To see locality under
reblocking, a bond meeting $X$ has, after division
by $L$, an endpoint in $Y$ whenever
$\overline X^{\,L}=LY$. Analyticity follows from
$C_L\in\mathfrak X_j$.

For $d_M(Y)\ge1$, the reblocking estimate gives
\[
 \sum_{\overline X^{\,L}=LY}\mathrm e^{-\kappa d_M(X)}
 \le c_{\rm rb}L^3\mathrm e^{-L(\kappa-\kappa_0-1)d_M(Y)}.
\]
With
$L(\kappa-\kappa_0-1)\ge\kappa$ and
$c_{\rm rb}L^3\mathrm e^{-[L(\kappa-\kappa_0-1)-\kappa]}\le1/2$,
the weighted supremum $q_j$ of $\mathfrak P_j$
over $d_M(X)\ge1$ satisfies
$q_{j+1}\le(q_j+2C_dM^3)/2$. Since $q_0=0$,
\begin{equation}\label{eq:interacting-P-bound}
 |\mathfrak P_j(X,C)|\le2C_dM^3\mathrm e^{-\kappa d_M(X)}
                           \qquad(d_M(X)\ge1).
\end{equation}
At $j=0$ the determinant activity includes all
$192$ chamber components. The primitive bound and
reblocking estimate have the same constants at
that first step.

The determinant identity
\cite[(33)]{NonlinearStep} gives
\begin{align}
 \frac{\det(D^{\rm ch}[qa]+\bar m\eta)}{\det(D^{\rm ch}[0]+\bar m\eta)}
 &=\exp\left(\sum_X\mathfrak P_*(X,a^{[k_*]})\right)
                    \frac{\det D_*(a^{[k_*]})}{\det D_*(0)},
                           \label{eq:interacting-determinant-telescope}\\
 \sum_X\mathfrak P_*(X,a^{[k_*]})
 &=\sum_{j<k_*}\sum_XE_j^{\rm det}(X,a^{[j]}).
                           \label{eq:interacting-P-telescope}
\end{align}
Here $a^{[j]}\in\mathfrak X_j$ for every $j$.
The fine determinant uses the phase-argument convention
of \eqref{eq:response-free-ratio}: its connection argument
is $qa(t)=\sqrt\eta\sum_ct_c\widehat{\mathfrak a}_c$.
For a single cell this is precisely the phase connection
$q(t/e)\widehat a_c$ in \eqref{eq:response-free-ratio}.
Write $n_j=\ell_j^3$ for the number of unit sites at level $j$;
the spinor dimension is $4n_j$ for $j\ge1$, and $192n_0$
at the incoming chamber level. The one-step identity reads
\[
 \det D_{j+1}(a^{[j+1]})
       =b^{4n_{j+1}}\det D_j(a^{[j]})/\delta Z_j(a^{[j]}).
\]
Divide it by its zero-field version and use
\cite[(78)]{FermionBlock}. The resulting
ratio acquires the factor
\[
 \exp\left(-\sum_XE_j^{\rm det}(X,a^{[j]})\right)
\]
at each
step, which proves the first identity by iteration.
Summing \eqref{eq:interacting-P-recursion} over $Y$
proves the second: each fine polymer has exactly
one reblocked image. The zero-field determinants
are nonzero by the real covariance formula.
At $j=0$ the incoming determinant has dimension
$192n_0$, the outgoing one dimension $4n_1$;
the factor $b^{4n_1}$ cancels in the ratio.

Fix $A\in V_{\mathrm t}(T_*^0)$ with $\chi_*(A)=1$
and a real twist $\theta$. In particular, $A$ is real and has mean zero.
The real gauge construction in
\cite[(32)]{MaxwellBlock}--%
\cite[(33)]{MaxwellBlock}
gives a representative $A^g$ with
$|A^g|_\infty\le2\ell_*^2p_*$.
The curvature bound and
\cite[Theorem~B]{MaxwellBlock}
put both $A+c_\theta$ and $A^g+c_\theta$
in $\mathrm{Dom}_{k_*}$.
Write
\[
 \mathscr A^0=H_*^{\rm ph}A^g,\qquad
 c_{\theta,\mu}=\frac{\theta_\mu}{e_*\ell_*},
 \qquad \mathscr A=\mathscr A^0+c_\theta.
\]
Summing the three photon kernel estimates in
\cite[(14)]{MaxwellBlock} bounds each
field norm of $\mathscr A^0$ by
$2C_0\ell_*^2p_*<e_*^{-\varepsilon}/2$;
the strict inequality is the corresponding
condition in \cite[(43)]{TerminalPressure}.
The local photon domains and their invariance
under real constants then give
$s\mathscr A^0+c_\theta\in\frac14\widetilde R_*$
for $0\le s\le1$. In particular all
connections needed below belong to $\mathcal W$.

The flow bound \cite[(66)]{CoverTuning}
gives an activity with even Grassmann kernels,
with equal barred and unbarred degrees, and
\begin{equation}\label{eq:interacting-activity-input}
 \|\mathscr E_*\|_*^x
       \le4e_*^{1/4-6\varepsilon},\qquad
 |m_*|\le e_*^{3/4-8\varepsilon}/8.
\end{equation}
The coefficient norm includes the weights
$\eta_*^3/48$ for a fine spinor index and
$(\eta_*^3/48)^2$ for a pair index, with all
admissible pairs, including pairs inside a cluster.
Its polymer factor is $\mathrm e^{\kappa d_M(X)}$.
For $(A,C)\in\widetilde R_*\times\mathfrak X_*$ define
\begin{equation}\label{eq:interacting-two-field-integral}
 Z_{\rm in}(A,C)=\int D\Psi\,
 \exp\left\{-\langle\bar\Psi,D_*(A+C)\Psi\rangle
       -m_*\langle\bar\psi,\psi\rangle
       +\sum_X\mathscr E_*(X;A,C,\psi,\chi)\right\},
\end{equation}
where $\psi,\bar\psi$ are substituted using
$H_*^\#(A+C)\Psi$ and their barred partners,
and $\chi,\bar\chi$ using the corresponding
covariant Hölder quotients.
The inclusion
$\widetilde R_*+\mathfrak X_*\subset\widetilde G_*$
makes all these maps analytic; the activity
kernels are analytic in $(A,C)$.
Berezin integration is coefficient extraction,
and the exponential is its scalar exponential
times a finite polynomial in the positive-degree
coefficients. Thus $Z_{\rm in}$ is analytic
on this two-field domain.
Set
\begin{align}
 Z_{\rm tm}(t)&=Z_{\rm in}(\mathscr A,\mathscr C(t)),
                               \label{eq:interacting-terminal-integral}\\
 f_{\rm int}(t)&=
   \frac{Z_{\rm tm}(t)}{\det D_*(c_\theta+\mathscr C(t))}.
                               \label{eq:interacting-factor-definition}
\end{align}
The Berezin measure has one pair for each
of the $4\ell_*^3$ coarse spinor coordinates.
The two-field analyticity gives analyticity of
$Z_{\rm tm}$ on
$\mathbb D(t_x^+)$.

\begin{proposition}\label{prop:interacting-factor}
Suppose
\begin{gather}
 4C_0\ell_*^3g_1\le e_*^{-1/4+\varepsilon},\qquad
 4\ell_*^3(2C_K+K_{\mathrm{pol}}+4)\le e_*^{-2\varepsilon},\qquad
 e_*^{1/4-9\varepsilon}\le1/8,
                                      \label{eq:interacting-factor-conditions}\\
 48e_*^{2\varepsilon}\le\mu_- .
                                      \label{eq:interacting-small-complex-mass}
\end{gather}
Then $f_{\rm int}$ is analytic and
\begin{equation}\label{eq:interacting-factor-bound}
 |f_{\rm int}(t)-1|\le\frac12e_*^{1/4-9\varepsilon}
                              \le\frac1{16}.
\end{equation}
In particular $Z_{\rm tm}$ has no zero on the
polydisc. The normalized terminal ratio
\begin{equation}\label{eq:interacting-terminal-ratio}
 R_{\rm tm}(t;\theta,A)=
 \frac{Y_*^{\rm x}(A+c_\theta;\mathscr C(t))}
      {Y_*^{\rm x}(A+c_\theta;0)}
 =\exp\left(\sum_X\mathfrak P_*(X,\mathscr C(t))\right)
       \frac{Z_{\rm tm}(t)}{Z_{\rm tm}(0)}
\end{equation}
is nonzero, continuous in $(t,\theta,A)$ on
$\mathbb D(t_x^+)\times\R^3\times\{A\in V_{\mathrm t}(T_*^0):\chi_*(A)=1\}$,
and analytic in $t$. It has a unique continuous
logarithm $\mathcal G(t;\theta,A)$ that vanishes
at zero, and this logarithm is analytic in $t$.
\end{proposition}

\begin{proof}
Use the covariance
\[
 C_t=D_*(c_\theta+\mathscr C(t))^{-1}.
\]
The inverse exists by
\cite[(75)]{TerminalPressure}, because
$e_*|\mathscr C(t)|_\infty<e_*^{3\varepsilon}
<\mu_*/24$. Each covariance row has sum at
most $4\ell_*^3g_1$. Put
$a_w=C_0^{-1}e_*^{-1/4+\varepsilon}$.
The first condition in
\eqref{eq:interacting-factor-conditions} gives
$a_w\ge4\ell_*^3g_1\ge1$.

Separate the quadratic form with covariance $C_t$
from the remaining even interaction $V_t$.
Normalized Gaussian Grassmann integration gives
\begin{equation}\label{eq:interacting-crude-Gaussian}
 |f_{\rm int}(t)-1|\le \mathrm e^{|V_t|_{a_w}}-1.
\end{equation}
To justify this particular norm bound, replace
the unbarred generators by $C_t\Psi$ and leave
the barred generators fixed. Their row sums
are at most $a_w$, so this substitution maps
the $a_w$ coefficient norm into the unit
coefficient norm. Integration against the identity
covariance has norm one in that latter norm.
Apply this to $\mathrm e^{V_t}-1$, and use the
submultiplicativity of the coefficient norm.

There are three terms in $V_t$. The difference
between the two coarse quadratic forms has
entry sum at most
$C_Ke_*^{1-\varepsilon}\ell_*^3/2$, by
\cite[(66)]{TerminalPressure} integrated along
$c_\theta+\mathscr C(t)+s\mathscr A^0$.
A quadratic Grassmann form has coefficient norm
at most $4a_w^2$ times this block sum, so
its contribution is bounded by
$2C_K\ell_*^3e_*^{1/2+\varepsilon}$.
The terminal mass estimate and
\eqref{eq:interacting-activity-input} give respectively
\[
 \frac12\ell_*^3e_*^{1/4-6\varepsilon},
 \qquad K_{\mathrm{pol}}\ell_*^3e_*^{1/4-7\varepsilon}.
\]
For the activity sum, substitute the terminal
field maps in the coefficient norm, then sum
$\mathrm e^{-\kappa d_M(X)}$ by assigning each polymer
one of its $M$-cubes and using the polymer
sum bound $K_{\mathrm{pol}}$. We have used
$4e_*^\varepsilon\le1$.
Consequently
\[
 |V_t|_{a_w}\le
 \ell_*^3(2C_K+K_{\mathrm{pol}}+4)e_*^{1/4-7\varepsilon}
 \le\frac14e_*^{1/4-9\varepsilon}.
\]
The inequality $\mathrm e^x-1\le2x$ proves
\eqref{eq:interacting-factor-bound}.
The denominator in \eqref{eq:interacting-factor-definition}
is analytic and nonzero, so $f_{\rm int}$ is analytic.

By \cite[(16)]{TerminalPressure} the terminal density has the form
\[
 Y_*^{\rm x}(A';C)=
 (\det D_*(0))^{-1}
 \exp\left(-\varepsilon_*^{\rm v}\ell_*^3+
                         \sum_X\mathfrak P_*(X,C)\right)
 Z_{\rm in}(H_*^{\rm ph}A',C).
\]
Its gauge invariance in \cite[Theorem~B]{TerminalPressure}
permits replacing $A$ by
$A^g$. The photon minimizer preserves real
constants. Dividing the density at $C=\mathscr C(t)$
by its value at zero proves
\eqref{eq:interacting-terminal-ratio}; the quotient
also shows that the result is independent of
the chosen gauge representative. For continuity,
use the ungauged linear map
$A+c_\theta\mapsto H_*^{\rm ph}(A+c_\theta)$
in the displayed formula. The terminal Grassmann
integral is analytic in its two field arguments,
and its zero-mark denominator is nonzero.

The required logarithm is
\begin{equation}\label{eq:interacting-logarithm}
 \mathcal G(t;\theta,A)=
 \int_0^1\sum_ct_c
       \frac{\partial_{t_c}R_{\rm tm}(st;\theta,A)}
            {R_{\rm tm}(st;\theta,A)}\,\dd s .
\end{equation}
Local Cauchy formulas show that the derivatives
in this integrand are jointly continuous;
integration over $[0,1]$ preserves continuity
and analyticity. Along a ray the integrand is
the logarithmic derivative of $R_{\rm tm}(st)$.
Differentiating the product of this ratio with
the exponential of minus the integrated
logarithmic derivative gives zero, and its
value at $s=0$ is one. Hence
$\mathrm e^{\mathcal G}=R_{\rm tm}$.
Any two continuous logarithms differ by
$2\pi i$ times an integer-valued continuous
function on the connected polydisc. Their
common value at zero proves uniqueness.
\end{proof}

\section{The one-cell logarithm and harmonic twists}
\label{subsec:interacting-one-cell}

For a fixed cell $c$ and real twist $\theta$, let
\begin{equation}\label{eq:interacting-free-coarse-log}
 g_D^\theta(t)=\tr\log\left(1+D_*(c_\theta)^{-1}
                   [D_*(c_\theta+t\widehat c_c)-D_*(c_\theta)]\right).
\end{equation}
The logarithm here is the matrix power series.
Set
\begin{align}
 C_{\rm fr}&=2592C_dK_{\mathrm{pol}}C_\beta^2
       \left(M^2(D_\beta^{\rm buf})^3L^4h^{-2}+8M^5h\right),
                                                   \notag\\
 B_{\rm wr}&=16\exp(6)C_dM^3K_{\mathrm{pol}}p_\beta \mathrm e^{-r_*},
 \qquad B_1=16g_1C_K'p_\beta e_*^{2\varepsilon},       \notag\\
 B_0&=C_{\rm fr}e_*^{4\varepsilon}+B_{\rm wr}+B_1 .
                                \label{eq:interacting-free-log-constants}
\end{align}

\begin{proposition}\label{prop:interacting-one-cell}
Suppose \eqref{eq:interacting-factor-conditions}--%
\eqref{eq:interacting-small-complex-mass} hold,
$r_*>31L+16$, and
\begin{equation}\label{eq:interacting-one-cell-conditions}
 2g_1C_K'e_*^{2\varepsilon}\le1/2,\quad
 B_0\le1/4,\quad
 \tau_s:=c_\tau\ell_*\ge1,\quad
 52B_1\mathrm e^{-\tau_s}\le e_*^{1/4-9\varepsilon}.
\end{equation}
The free ratio $N_0$ has the logarithm
\begin{equation}\label{eq:interacting-one-cell-free-log}
 \ell_0(t)=\sum_X\mathfrak P_*(X,t\widehat c_c)+g_D^0(t),
 \qquad N_0(t)=\mathrm e^{\ell_0(t)} ,
\end{equation}
analytic on $|t|<t_x^+$. It is independent of
$c$ and of the admitted placement, is even
and real on the real axis, and satisfies
\begin{align}
 \left|\sum_X\mathfrak P_*(X,t\widehat c_c)\right|
        &\le C_{\rm fr}e_*^{-2\varepsilon}|t|^2+B_{\rm wr},
                                  \label{eq:interacting-P-probe-gain}\\
 |\ell_0(t)|&\le B_0\qquad(|t|\le t_x).
                                  \label{eq:interacting-free-log-bound}
\end{align}
The one-cell interacting error
$w_c(t)=\mathcal G(te_c;\theta,A)-\ell_0(t)$
is analytic, vanishes at zero, and obeys
\begin{equation}\label{eq:interacting-one-cell-error}
 |w_c(t)|\le3e_*^{1/4-9\varepsilon}
                         \qquad(|t|<t_x^+).
\end{equation}
\end{proposition}

\begin{proof}
Write $K=D_*(c_\theta+t\widehat c_c)-D_*(c_\theta)$.
The real covariance and the integrated insertion
bounds give
\[
 \|D_*(c_\theta)^{-1}K\|_{\mathrm{op}}\le
           2g_1C_Ke_*^{2\varepsilon}\le1/2,\qquad
 \|D_*(c_\theta)^{-1}K\|_{\mathfrak S_1}
           \le4g_1C_Kp_\beta e_*^{2\varepsilon}.
\]
Thus the series in \eqref{eq:interacting-free-coarse-log}
converges, is analytic, vanishes at zero, and
has modulus at most $B_1$. Triangularizing
the matrix identifies its exponential with
$\det D_*(c_\theta+t\widehat c_c)/\det D_*(c_\theta)$.
Combine this with
\eqref{eq:interacting-determinant-telescope}.
The lattice symmetries carrying one cell
probe to another preserve the fine determinant,
so the resulting fine ratio is $N_0$.

The bound \eqref{eq:interacting-P-probe-gain}
requires a cancellation before summing the levels.
Put $i=k_*-j$ and $B_j=e^{-1}\widehat{\mathfrak a}_c^{[j]}$.
Rescaling gives
\begin{equation}\label{eq:interacting-probe-derivative}
 e_j|\partial B_j|_\infty\le
                         a_j^2a_*^{-3/2}C_\beta .
\end{equation}
Let $\mathcal Q_j$ be the $M$-cubes containing
endpoints of the supporting bonds of $B_j$.
The buffered support cube yields
\begin{equation}\label{eq:interacting-probe-cube-count}
 |\mathcal Q_j|\le
 4\left((D_\beta^{\rm buf})^3M^{-3}a_j^{-3}+8\right),
 \qquad |\mathcal Q_j|\le L^{3i}p_\beta .
\end{equation}
By locality only polymers containing a cube
of $\mathcal Q_j$ contribute to
\eqref{eq:interacting-P-telescope}.

For such a polymer $X$, put $d=d_M(X)$ and
$s_X=M(d+3)$. Suppose first that $2s_X<\ell_j$.
All bonds meeting $X$ lie in a coordinate
cube of radius $M(d+2)+1$ around a chosen
site $\tilde c$ of $X$. The strict inequality
$2[M(d+2)+1]+\eta_j<2s_X<\ell_j$
makes this cube a lift to the lattice.
Define a real function there by
$\lambda(c)=\sum_\nu B_{j,\nu}(\tilde c)(c-\tilde c)_\nu$,
and extend it by zero outside. On bonds meeting
$X$, its gradient is the constant value
$B_j(\tilde c)$. If $[\,\cdot\,]_X$ restricts
a connection to those bonds, then
\[
 B_X=[B_j-d\lambda]_X,\qquad
 |B_X|_\infty\le3s_X|\partial B_j|_\infty.
\]
The last inequality follows by integrating
the first differences along coordinate paths
inside the lifted cube.

Real gauge invariance must be used before
continuing to the complex gauge parameter.
On the convex open set where
$|[tB_j-zd\lambda]_X|_\infty<e_j^{-1+\varepsilon}$,
the analytic function
$z\mapsto E_j^d(X,[tB_j-zd\lambda]_X)$
equals $E_j^d(X,tB_j)$ for real $z$ near zero:
the unrestricted connection is related to
$tB_j$ by a real gauge transformation, and
locality permits the restriction. The identity
theorem gives this equality throughout that
set. In particular it holds at $z=t$
whenever $|t|<\rho_X$, where
$\rho_X=e_j^{-1+\varepsilon}/(3s_X|\partial B_j|_\infty)$.

The function
$z\mapsto E_j^d(X,zB_X)-E_j^d(X,0)$ is even,
bounded by $2C_dM^3\mathrm e^{-(\kappa+1)d}$ on
$|z|<\rho_X$, and has a double zero at zero.
Apply the maximum principle after dividing
by $z^2$. For $|t|\ge\rho_X$ the original
primitive bound gives the same inequality
with the factor $(|t|/\rho_X)^2\ge1$.
Thus, in both cases,
\begin{equation}\label{eq:interacting-local-constant-removal}
 |E_j^{\rm det}(X,tB_j)|\le
 18C_dM^5C_\beta^2(d+3)^2\mathrm e^{-(\kappa+1)d}
          |t|^2e_j^{-2\varepsilon}a_j^4a_*^{-3}.
\end{equation}
When $\partial B_j=0$, $B_X=0$ and this contribution
vanishes. This proves the cancellation also
at the chamber first step.

Since $(d+3)^2\mathrm e^{-(\kappa+1)d}\le9\mathrm e^{-\kappa_0d}$,
summing \eqref{eq:interacting-local-constant-removal}
over cubes and polymers gives at level $j$
the bound
\[
 648C_dK_{\mathrm{pol}}C_\beta^2|t|^2e_j^{-2\varepsilon}a_*^{-3}
           \left(M^2(D_\beta^{\rm buf})^3a_j+8M^5a_j^4\right).
\]
For the remaining polymers $2s_X\ge\ell_j$,
$d\ge r_*L^i/2-3$. The extra exponential in
the primitive bound then gives at level $j$
\[
 2\exp(6)C_dM^3K_{\mathrm{pol}}p_\beta L^{3i}\mathrm e^{-r_*L^i}.
\]
Sum over $i\ge1$, using $a_j=L^{-i}a_*$,
$e_j^{-2\varepsilon}=e_*^{-2\varepsilon}L^{i\varepsilon}$,
$\sum_iL^{-i(1-\varepsilon)}\le1$, and
$a_*^{-2}\le L^4h^{-2}$, $a_*\le h$.
For $r_*>31L+16$, the second sum is bounded
by a constant times $\mathrm e^{-r_*}$.
The stated definitions of $C_{\rm fr}$ and
$B_{\rm wr}$ give \eqref{eq:interacting-P-probe-gain}.
Together with $|g_D^0|\le B_1$ this proves
\eqref{eq:interacting-free-log-bound}.

Evenness and reality follow from the corresponding
fine determinant symmetries. Both $\ell_0(t)$
and $\ell_0(-t)$ vanish at zero and have
exponential $N_0(t)$, so they coincide.
For real $t$, the nonzero real ratio $N_0(t)$
is positive by continuity from one at zero;
the normalized logarithm is therefore real.
Cell and placement independence follows in
the same way from the determinant translation
and reflection symmetries.

It remains to control the twist. For
$|\operatorname{Im}\theta_\mu|\le\tau_s$
the connections $c_\theta+st\widehat c_c$
lie in $\mathcal W^+$, and
$\|D_*(c_\theta)^{-1}\|_{\mathrm{op}}\le g_1$.
There is a positive strip margin
$\ell_*(e_*^{2\varepsilon}-e_*^{3\varepsilon})$:
split the imaginary constant connection
into a part bounded by $c_\tau/e_*$ and a
remaining part absorbed into $\mathfrak X_*$.
The latter has vanishing first differences.
The mass condition bounds its total imaginary
part by $\mu_*/24$. On this larger open strip,
\cite[(73)]{TerminalPressure} gives
the same matrix power-series bound $|g_D^\theta|\le B_1$.
The function is holomorphic there and
$2\pi$-periodic in each twist variable.
For real twists, periodicity is conjugation by
the periodic gauge factor
$\exp(-2\pi i c_\mu/\ell_*)$; the identity
theorem extends it to the strip.

We recall the short Fourier estimate to make the
size of this suppression explicit. A bounded
periodic holomorphic function $g$ in a closed
strip of width $\tau\ge1$ has
$|\widehat g(m)|\le B \mathrm e^{-\tau|m|_1}$:
shift each Fourier contour towards
$-i\tau\operatorname{sgn}(m_\mu)$, first using
an interior width and then taking its limit.
The vertical sides cancel by periodicity.
The nonconstant Fourier sum is bounded by
\[
 B\left[\left(\frac{1+\mathrm e^{-\tau}}{1-\mathrm e^{-\tau}}\right)^3-1\right]
                                      \le25B\mathrm e^{-\tau}.
\]
The last inequality follows by expanding the
bracket divided by $\mathrm e^{-\tau}$ as a power
series with positive coefficients and evaluating
at $\mathrm e^{-\tau}=\mathrm e^{-1}$. Thus the difference
between two real values is at most $52B\mathrm e^{-\tau}$.
Apply this to $g_D^\theta(t)$ to obtain
\begin{equation}\label{eq:interacting-twist-suppression}
 |g_D^\theta(t)-g_D^0(t)|\le52B_1\mathrm e^{-\tau_s}
                                  \le e_*^{1/4-9\varepsilon}.
\end{equation}

Finally the normalized logarithm identity and
\eqref{eq:interacting-factor-definition} give
\[
 w_c(t)=g_D^\theta(t)-g_D^0(t)
                +\log f_{\rm int}(te_c)-\log f_{\rm int}(0).
\]
The principal logarithm of $f_{\rm int}$ is
analytic and has modulus at most
$e_*^{1/4-9\varepsilon}$. Both sides vanish
at zero and exponentiate to the same ratio;
uniqueness of the logarithm justifies the
identity. Equation
\eqref{eq:interacting-twist-suppression} proves
\eqref{eq:interacting-one-cell-error}.
\end{proof}

Let $\phi_0(t)=\sum_{i=0}^4N_0^{(i)}(0)t^i/i!$.
Cauchy's estimate applied to
$|N_0(t)-1|\le2B_0$ on $|t|=t_x=4r$ gives
\begin{equation}\label{eq:interacting-free-jet}
 |c_2|\le r^{-2}/16,\qquad |c_4|\le r^{-4}/16,\qquad
 \|\phi_0-1\|_r\le1/6,\qquad \|\phi_0\|_r\le2.
\end{equation}
The odd coefficients vanish. Expanding $\log N_0$
through degree four gives
$\ell_0^{(4)}(0)=c_4-3c_2^2$, the chamber fourth coefficient
of \cite[Theorem~B]{PositiveProbe}.

\section{Derivatives in separated cells}
\label{subsec:interacting-separated}

To estimate mixed logarithmic derivatives, we use the
real covariance $G=D_*(\mathscr A)^{-1}$. On the
coarse Grassmann algebra, write an even element
without constant term as
\[
 F=\sum_{p\ge2}\frac1{p!}\sum_{u_1,\ldots,u_p}
 F_p(u_1,\ldots,u_p)\Psi(u_1)\cdots\Psi(u_p),
\]
with antisymmetric kernels.
For a symmetric weight $w\ge1$, define
\begin{align}
 \rho_G(F;w)&=\sum_{p\ge2}\frac{h_G^p}{(p-1)!}
       \sup_{u_1}\sum_{u_2,\ldots,u_p}|F_p(u_1,\ldots,u_p)|w,
                                                     \notag\\
 N_G(F;w)&=\sum_{p\ge2}\frac{h_G^p}{(p-1)!}
       \sum_{u_1,\ldots,u_p}|F_p(u_1,\ldots,u_p)|w .
                                     \label{eq:interacting-tree-norms}
\end{align}
The supremum in $\rho_G$ is taken separately
at each degree. The indices include site, spin,
and barred or unbarred status; their distance
is the distance between the sites. We use
\[
 w_\delta(u_1,\ldots,u_p)=
          \mathrm e^{\vartheta_0\max_{i,j}d(u_i,u_j)},\qquad
 w_P(u_1,\ldots,u_p)=\mathrm e^{\vartheta_0\max_i d(u_i,P)} .
\]
The constants $h_G,\vartheta_0,\alpha_G$ are
those of \cite[(74)]{TerminalPressure}.
Antisymmetrizing a kernel that is not antisymmetric
bounds $N_G$ by the same expression for that
kernel. For $\rho_G$, replace the supremum at
the first argument by the largest supremum
over all argument positions.

In the fine algebra absorb the terminal mass
into each single-cube activity:
\[
 E^a(X;t)=\mathscr E_*(X;\mathscr A,\mathscr C(t))
              -\mathbf1_{\{X\text{ is an }M\text{-cube}\}}
                         m_*\langle\bar\psi,\psi\rangle_X.
\]
Let $E^a_{nm}(X;t)$ denote its fine spinor/pair
kernels. The antisymmetric mass kernel has
norm $8M^3$: its two fine-index weights multiply
$48\eta_*^{-3}$, its diagonal Kronecker factors
leave $48(M/\eta_*)^3$ chamber sites, and there
are four spin components and two orders.
It follows from \eqref{eq:interacting-activity-input}
that, when $a_w\le e_*^{-1/4+\varepsilon}$ and
$(4+M^3/2)e_*^\varepsilon\le1$,
\begin{equation}\label{eq:interacting-absorbed-activity}
 \sum_{n,m}\frac{a_w^{n+m}}{n!m!}\|E^a_{nm}(X;t)\|
          \le w_*^{\rm E}\mathrm e^{-\kappa d_M(X)},\qquad
 w_*^{\rm E}=e_*^{1/4-7\varepsilon}
                                      \quad(|t|_\infty\le t_x).
\end{equation}
Cauchy's formula on the circles of radius $t_x/2$
in one or two distinct mark variables gives,
on $|t|_\infty\le t_x/2$, the same bound
for $\partial_cE^a_{nm}$ and
$\partial_c\partial_{c'}E^a_{nm}$ multiplied
respectively by $q_x=2/t_x$ and $q_x^2$.
The circle integrals precede the finite weighted
kernel sums, so the coefficient norm is bounded
without replacing it by a supremum of individual
coefficients. Locality gives
$\partial_cE^a(X;t)=0$ unless $X\cap P_c\ne\varnothing$.

After substituting the terminal field maps, put
\begin{align}
 V(t)&=-\langle\bar\Psi,K(t)\Psi\rangle
           +\sum_X E^a(X;t)[H_*^\#(\mathscr A+\mathscr C(t))\Psi],
                                                        \notag\\
 K(t)&=D_*(\mathscr A+\mathscr C(t))-D_*(\mathscr A),\qquad
 V_0(t)=\sum_XE^a_{00}(X;t),\qquad V_+(t)=V(t)-V_0(t).
                                    \label{eq:interacting-tree-interaction}
\end{align}
It is even and analytic, and $V_+$ has no
constant term. Gaussian integration gives
\begin{equation}\label{eq:interacting-real-Gaussian-factor}
 Z_{\rm tm}(t)=\det D_*(\mathscr A)\mathrm e^{V_0(t)}
                                  \int\dd\mu_G\,\mathrm e^{V_+(t)} .
\end{equation}
The choice of the real covariance will allow
the exponential weights in the tree expansion.

\begin{proposition}\label{prop:interacting-differentiated-norms}
Suppose
\begin{equation}\label{eq:interacting-tree-smallness}
 5h_{\rm L}h_G\le e_*^{-1/4+\varepsilon},\qquad
 (4+M^3/2)e_*^\varepsilon\le1,\qquad
 \alpha_G\bar r\le1/2,
 \quad
 \bar r=4h_G^2C_Ke_*^{2\varepsilon}
                            +650K_{\mathrm{pol}}e_*^{1/4-7\varepsilon}.
\end{equation}
Let $d_{\rm sep}=\min_{c\ne c'}d(P_c,P_{c'})$ and
suppose $d_{\rm sep}\ge6M+8$.
With
$s_F=(24/\gamma^F)^3$ and
$C_V^{\mathrm{ins}}=10^3C_H^2s_F^4K_{\mathrm{pol}}$, we have
\begin{align}
 \rho_G(V_+(t);w_\delta)&\le\bar r
                                      &&(|t|_\infty\le t_x),\notag\\
 N_G(\partial_cV_+(t);w_{P_c})&\le q_xC_V^{\mathrm{ins}}h_G^2p_\beta,
 \quad
 \rho_G(\partial_cV_+(t);w_{P_c})\le q_xC_V^{\mathrm{ins}}h_G^2,\notag\\
 N_G(\partial_c\partial_{c'}V_+(t);1)
          +|\partial_c\partial_{c'}V_0(t)|
 &\le q_x^2C_V^{\mathrm{ins}}h_G^2p_\beta
                    \mathrm e^{-\vartheta_0(d_{\rm sep}-4M-6)} .
                             \label{eq:interacting-differentiated-norms}
\end{align}
The last two lines hold on $|t|_\infty\le t_x/2$,
with $c\ne c'$ in the last line. If the lattice
probe vanishes, all marked derivatives vanish.
\end{proposition}

\begin{proof}
We give the kernel summation which preserves
the marked supports. For a polymer $X$, let
$X^+$ be its unit-site one-neighborhood.
Every fine index in an activity supported on
$X$ has its position set in $X^+$. This
includes both chamber sites in an internal
pair and the first site of an external
nearest-neighbor pair whose other endpoint
lies over $X$. The polymer diameter estimate
gives
\begin{equation}\label{eq:interacting-X-diameter}
 \vartheta_0\operatorname{diam}(X^+)\le d_M(X)+3 .
\end{equation}
The weighted kernel after substituting the
field maps is, at degree $p=n+m$,
\[
 K_{nm}^X(u_1,\ldots,u_p)=
 \int E^a_{nm}(X;t;v_1,\ldots,v_p)
           \prod_{i=1}^p\operatorname{Leg}(v_i;u_i)\,\dd v .
\]
Here $\dd v$ includes the fine weights already
specified. Passing from the bidegree expression
to a degree-$p$ kernel produces the factor
$p!/(n!m!)$; in the tree norms its remaining
factor is $p h_G^p/(n!m!)$.

Suppose a weight is bounded by
$c_w\prod_i \mathrm e^{\vartheta_0d(u_i,\operatorname{pos}(v_i))}$.
Each summed leg contributes at most $h_{\rm L}$.
If one argument $u$ is fixed, its leg contributes
at most
$h_{\rm L}\mathrm e^{-(\gamma^F/4)d(u,X^+)}$.
Thus the full kernel sum is at most
$c_wh_{\rm L}^p\|E^a_{nm}(X;t)\|$, and the
sum with one fixed argument retains the
extra exponential in $d(u,X^+)$.
For $w_\delta$, the triangle inequality between
the two farthest arguments and their position
sets, followed by \eqref{eq:interacting-X-diameter},
permits $c_w=\mathrm e^{d_M(X)+3}$.

The needed rooted polymer sum is
\begin{equation}\label{eq:interacting-rooted-polymer-sum}
 \sum_X \mathrm e^{-(\gamma^F/4)d(u,X^+)}
                       \mathrm e^{-(\kappa-1)d_M(X)}\le32K_{\mathrm{pol}}.
\end{equation}
Indeed the sum of
$\mathrm e^{-(\gamma^F/4)d(u,\square^+)}$ over $M$-cubes
is at most $32$. The neighboring $27$ cubes
contribute at most $27$; a cube $j$ layers
away has distance at least $(j-1)M$ and there
are at most $(2j+1)^3$ such cubes.
Use $\gamma^FM/4\ge4$ to sum the remaining
geometric tail. Assigning a cube of $X$
and applying the anchored polymer sum $K_{\mathrm{pol}}$
proves \eqref{eq:interacting-rooted-polymer-sum}.

Apply \eqref{eq:interacting-absorbed-activity}
with $a_w=5h_{\rm L}h_G$. The factor $5^{-p}$
absorbs the polynomial degree factors.
The preceding full and rooted sums give
\[
 \rho_G\left(\sum_X E^a(X;t)[H_*^\#\Psi]_+;w_\delta\right)
                                       \le650K_{\mathrm{pol}}w_*^{\rm E}.
\]
More generally, for kernels bounded by
$\mathcal B \mathrm e^{-\kappa d_M(X)}$ at that
weight, these sums give
$650K_{\mathrm{pol}}\mathcal B$ in $\rho_G$ and
$\exp(3)\mathcal B\sum_X\mathrm e^{-(\kappa-1)d_M(X)}$
in $N_G$. This follows by summing
$\sum_{p\ge2}p5^{-p}=9/80$ in the rooted
case and using $p\le5^p$ in the full case.
The quadratic term contributes at most
$4h_G^2C_Ke_*^{2\varepsilon}$, since
$\vartheta_0\le\gamma^F/2$. This proves the
first line of \eqref{eq:interacting-differentiated-norms}.

For the first derivative, there are three
types of term: the coarse quadratic derivative,
the activity-kernel derivative, and a derivative
of one leg. The first has respective bounds
$8h_G^2C_K\epsilon_cp_\beta$ and
$4h_G^2C_K\epsilon_c$ in $N_G$ and $\rho_G$,
with $\epsilon_c=e_*|\widehat c_c|_\infty\le q_x$.
In its chain through $z\in P_c$, the weight
$w_{P_c}$ is absorbed by the exponential
distances from the two coarse endpoints to $z$.

For the kernel derivative, $X$ must meet $P_c$.
The same weight comparison uses that meeting
point and \eqref{eq:interacting-X-diameter}.
The Cauchy bound multiplies
$\mathcal B=w_*^{\rm E}$ by $q_x$, and
the sum of polymers meeting $P_c$ is at most
$K_{\mathrm{pol}}p_\beta$. Its $N_G$ bound is
$\exp(3)K_{\mathrm{pol}}p_\beta q_xw_*^{\rm E}$ and its
$\rho_G$ bound is $650K_{\mathrm{pol}}q_xw_*^{\rm E}$.

For a differentiated leg at position $j$, the
insertion estimate has a chain
$w\in\operatorname{pos}(v_j)\to z\in P_c^+\to u_j$.
The weight obeys
\[
 w_{P_c}\le \mathrm e^{d_M(X)+5}
 \mathrm e^{\vartheta_0[d(w,z)+d(z,u_j)]}
            \prod_{i\ne j}\mathrm e^{\vartheta_0d(u_i,\operatorname{pos}(v_i))}.
\]
The extra two in the exponent accounts for
$d(P_c^+,P_c)\le M+1$.
Summing the differentiated leg and the other
legs leaves, before the degree sum,
\[
 2\mathrm e^{d_M(X)+5}h_{\rm L}^{p-1}
       C_H\epsilon_cs_F
       \sum_{z\in P_c^+}\mathrm e^{-(\gamma^F/2)d(z,X^+)}
                                      \|E^a_{nm}(X;t)\|.
\]
There are $p$ choices of the differentiated
leg. The kernel norm contributes one further
factor $p$, and $p^2\le5^p$ absorbs both.
Use \eqref{eq:interacting-rooted-polymer-sum}
with $\gamma^F/2$ in place of $\gamma^F/4$.
The result is bounded by
$9500C_Hs_FK_{\mathrm{pol}}w_*^{\rm E}p_\beta\epsilon_c$
in $N_G$.

For $\rho_G$, if the fixed argument lies on
the differentiated leg, split its two-edge
exponent as $7\gamma^F/8=3\gamma^F/8+\gamma^F/2$.
Summing the intermediate site leaves
$s_F\mathrm e^{-(\gamma^F/4)d(u,X^+)}$.
If the fixed argument lies on another leg,
that leg supplies the same rooted exponential,
and the differentiated leg contributes
$2C_H\epsilon_cs_F^2$.
Both cases therefore have the bound
$9500C_Hs_F^2K_{\mathrm{pol}}w_*^{\rm E}\epsilon_c$.
Here the degree sum is
$\sum_{p\ge2}p^25^{-p}=43/160$.
Adding the three types of term, using
$w_*^{\rm E}\le1/8$, $C_K\le C_Hs_F^2$,
and $C_H\ge2$, proves the second line with
$C_V^{\mathrm{ins}}=10^3C_H^2s_F^4K_{\mathrm{pol}}$.

For the mixed derivative, put
\[
 d_0=d(P_c^+,P_{c'}^+),\quad
 d_1=d(X^+,P_c^+),\quad d_2=d(X^+,P_{c'}^+).
\]
Then $d_0\ge d_{\rm sep}-2M-2$.
The differentiated coarse quadratic term is
bounded in $N_G$ by
\[
 16h_G^2C_H^2s_F^3p_\beta q_x^2
                                      \mathrm e^{-(\gamma^F/2)d_{\rm sep}}.
\]
For an activity kernel with $p$ legs, there
are $(p+1)^2$ placements of the two derivatives.
They may act on two legs, the same leg, a
kernel and a leg, or both on the kernel.
The summed first differentiated legs are
bounded by $2C_H\epsilon_cs_F^2\mathrm e^{-(\gamma^F/2)d_1}$
and its $c'$ version. The mixed leg is
bounded by
\[
 2C_H^2\epsilon_c\epsilon_{c'}s_F^3\mathrm e^{-(\gamma^F/2)d_0}
       \bigl(\mathrm e^{-(\gamma^F/2)d_1}+\mathrm e^{-(\gamma^F/2)d_2}\bigr).
\]
If a derivative acts on the kernel, its
corresponding distance $d_i$ is zero by
locality. Normalize the first kernel
derivatives by $q_x$ and the mixed one by
$q_x^2$. Each of the four resulting kernel
types satisfies \eqref{eq:interacting-absorbed-activity}.
Every product-rule term is consequently at most
$4q_x^2C_H^2s_F^4h_{\rm L}^p\|E_{\rm type}\|\mathcal D(X)$,
where
\[
 \mathcal D(X)=\mathrm e^{-(\gamma^F/2)(d_1+d_2)}
       +\mathrm e^{-(\gamma^F/2)d_0}
                  \bigl(\mathrm e^{-(\gamma^F/2)d_1}+\mathrm e^{-(\gamma^F/2)d_2}\bigr).
\]

The triangle inequality through $X^+$ gives
$d_0\le d_1+d_2+M d_M(X)+2M+2$.
Use $\vartheta_0\le\gamma^F/8$ and
$\vartheta_0M\le1$ to obtain
\[
 \mathcal D(X)\mathrm e^{-\kappa d_M(X)}
 \le2\mathrm e^{-\vartheta_0(d_0-2M-2)}
       \bigl(\mathrm e^{-(\gamma^F/4)d_1}+\mathrm e^{-(\gamma^F/4)d_2}\bigr)
                                \mathrm e^{-(\kappa-1)d_M(X)}.
\]
Sum the two distance factors over the sites
of $P_c^+$ and $P_{c'}^+$, then use
\eqref{eq:interacting-rooted-polymer-sum}.
The degree factor $p(p+1)^2$ is at most
$5^p$ for $p\ge2$. At $p=0$ there is
one kernel term with no legs, and the same
bound controls $\partial_c\partial_{c'}V_0$.
The four kernel types contribute in total
at most
$256C_H^2s_F^4K_{\mathrm{pol}}p_\beta q_x^2
\mathrm e^{-\vartheta_0(d_{\rm sep}-4M-4)}$.
Adding the quadratic term proves the last
line of \eqref{eq:interacting-differentiated-norms}.
If the lattice probe is zero, $\mathscr C(t)=0$
and there is no mark dependence.
\end{proof}

We next turn these kernel estimates into a bound
for the logarithm. The finite-dimensional tree
formula is recalled with its proof, since the
distance between its two marked vertices is
the source of the mixed-cell decay.

For $n$ copies of the generators, let
$\Delta^{ij}=\sum_{x,y}G(x,y)
\partial_{\bar\Psi^{(j)}(y)}\partial_{\Psi^{(i)}(x)}$
and $\Delta_{\{i,j\}}=\Delta^{ij}+\Delta^{ji}$.
The derivatives are left Grassmann derivatives.
For a forest $T$ and edge parameters
$\tau\in[0,1]^T$, let $s_{ij}^T(\tau)$ be the
minimum parameter on the path from $i$ to $j$,
zero if there is no path, and $s_{ii}^T=1$.
Write $\Delta[s]=\sum_{i,j}s_{ij}\Delta^{ij}$
and let $\epsilon_0$ extract the constant term.
For even elements $F_1,\ldots,F_n$ without
constant term, define
\begin{equation}\label{eq:interacting-connected-tree}
 T_n(F_1,\ldots,F_n)=
 \sum_{\text{trees }T}\int_{[0,1]^T}
 \epsilon_0\left\{\prod_{\ell\in T}\Delta_\ell
          \mathrm e^{\Delta[s^T(\tau)]}\prod_iF_i(\Psi^{(i)})\right\}\,\dd\tau .
\end{equation}
For $n=1$ this is the Gaussian expectation.

\begin{proposition}\label{prop:interacting-tree-log}
If $V_+(t)$ satisfies
\eqref{eq:interacting-tree-smallness} and is analytic,
then
\[
 W(t)=V_0(t)+\sum_{n\ge1}\frac1{n!}
                       T_n(V_+(t),\ldots,V_+(t))
\]
is analytic on $\mathbb D(t_x)$ and
$Z_{\rm tm}(t)=\det D_*(\mathscr A)\mathrm e^{W(t)}$.
For $c\ne c'$ and nonempty unit-site sets $P,P'$,
\begin{align}
 |\partial_c\partial_{c'}(W-V_0)|
 \le{}&2N_G(\partial_c\partial_{c'}V_+;1)\notag\\
 &+2\alpha_GN_G(\partial_cV_+;w_P)
       \rho_G(\partial_{c'}V_+;w_{P'})
                                  \mathrm e^{-\vartheta_0d(P,P')}.
                                  \label{eq:interacting-tree-mixed}
\end{align}
\end{proposition}

\begin{proof}
The forest interpolation formula for a polynomial $H$ is
\[
 H(\mathbf1)=\sum_F\int
       \left(\prod_{\ell\in F}\partial_{s_\ell}H\right)
                                 (s^F(\tau))\,\dd\tau.
\]
For a partition $\pi$ of the vertices, keep fixed
entries $s_{\mathrm{in}}(i,j)$ inside its blocks and
put $X_\pi(u)_{ij}=u$ on pairs joining distinct blocks.
The chain rule and the fundamental theorem of calculus give
\begin{equation}\label{eq:interacting-forest-first-step}
 G(X_\pi(u))=G(X_\pi(0))
       +\sum_{\substack{\ell=(i,j)\\
                        i,j\text{ in distinct blocks of }\pi}}
          \int_0^u(\partial_{s_\ell}G)(X_\pi(w))\,\dd w .
\end{equation}
When $\ell$ joins two blocks, merge them and assign
the value $w$ to every pair between those two blocks.
The interpolation point is unchanged. Begin with singleton
blocks, $G=H$ and $u=1$, and iterate this identity
on each integrand. Each derivative
merges two components, so after at most $n-1$
steps the procedure ends. The successive
parameters are ordered decreasingly; when two
vertices merge, their common value is the
minimum parameter on their eventual path.
For a fixed forest, its edge orderings divide
the parameter cube into these ordered simplices.
Summing them gives the formula.

The matrix $s^F$ is a Gram matrix of unit
vectors: at $u\in[0,1]$, assign to vertex $i$
the basis vector of its component among the
edges with $\tau_\ell>u$, and integrate their
scalar products in $u$. This scalar product
is the length of the interval on which the
two vertices remain connected, namely
$s_{ij}^F$. Wick's formula and the Gram
inequality therefore bound a residual monomial
of degree $p$ under $\mathrm e^{\Delta[s^F]}$
by $g^p$, where $\|G\|_{\mathrm{op}}\le g^2$.
Indeed, its determinant entries are
$\langle\omega_i\otimes\delta_x,\omega_j\otimes G\delta_y\rangle$.
These facts apply whether or not the
interpolated covariance is invertible.

Let vertex $i$ of a tree have degree $k_i$
and let the chosen homogeneous part of $F_i$
have degree $p_i$. The edge derivatives leave
the factorial $(p_i-k_i)!^{-1}$ and $p_i-k_i$
uncontracted fields at that vertex; terms with
$p_i<k_i$ vanish. Wick's Gram bound supplies
$g^{p_i-k_i}$. Sum contracted arguments from
the leaves to vertex $1$, using the weighted
row and column bound $\alpha_G$ for every
edge. If nonnegative symmetric weights $w_i$
and a distance $D$ satisfy
\[
 \mathrm e^{\vartheta_0D}\le
       \prod_iw_i(\text{arguments at }i)
       \prod_{\ell=(i,j)}
          \mathrm e^{\vartheta_0d(\text{edge arguments})},
\]
this yields $\mathrm e^{-\vartheta_0D}\alpha_G^{n-1}$
times the full kernel sum at vertex $1$
and the rooted sums at the other vertices.
The number of trees with degrees $k_1,\ldots,k_n$
is $(n-2)!/\prod_i(k_i-1)!$; here $n\ge2$,
$k_i\ge1$ and $\sum_i k_i=2n-2$, as in
\cite[Proposition~4]{PostnikovLin2019}.
Enlarge the sum over degrees and use
\[
 \sum_{k=1}^{p}
       \frac{g^{p-k}}{(k-1)!(p-k)!}
       =\frac{(1+g)^{p-1}}{(p-1)!}
       \le\frac{h_G^p}{(p-1)!}.
\]
We obtain, for $n\ge2$,
\begin{equation}\label{eq:interacting-tree-bound}
 |T_n(F_1,\ldots,F_n)|
 \le \mathrm e^{-\vartheta_0D}(n-2)!\alpha_G^{n-1}
                N_G(F_1;w_1)\prod_{i=2}^n\rho_G(F_i;w_i).
\end{equation}
For $n=1$ the bound is $N_G(F_1;1)$.
This also explains why the tree norms use
$(p-1)!$ rather than $p!$.

Apply the forest formula to the replicated
Gaussian expectation of $V_+^n$.
Different forest components contain disjoint
generator sets and their even factors commute,
so their expectations factor. Grouping by
component partitions gives
\[
 \int\dd\mu_G\,V_+^n
       =\sum_{\{B_1,\ldots,B_q\}}
                     \prod_{a=1}^qT_{|B_a|}(V_+,\ldots,V_+).
\]
The exponential formula for partitions
identifies the Gaussian integral with the
exponential of the connected tree sum.
The identity first holds as a formal series
in a scalar multiplying $V_+$. Bound
\eqref{eq:interacting-tree-bound} and
$\alpha_G\rho_G(V_+;1)\le1/2$ make the
connected series absolutely convergent up
to that scalar equal to one, so the formal
identity is an analytic identity there.
The tree terms are polynomial in the
activity coefficients, and their series
converges locally uniformly in $t$. Thus $W$
is analytic and its derivatives are obtained
term by term.

Two derivatives of the symmetric multilinear
tree expression either act on the same vertex
or on two different vertices. The first
terms, including $n=1$, sum to at most
$2N_G(\partial_c\partial_{c'}V_+;1)$.
For the second terms put the two derivatives
at vertices $1$ and $2$. The weights there
are $w_P,w_{P'}$, and those at all other
vertices are $w_\delta$.
Along the unique path from $1$ to $2$,
the distance $d(P,P')$ is at most the
sum of the two endpoint distances, all
edge distances, and the argument diameters
at intermediate vertices. Hence the weight
condition for \eqref{eq:interacting-tree-bound}
holds with $D=d(P,P')$. The tree sum leaves
$\alpha_GN_G(\partial_cV_+;w_P)
\rho_G(\partial_{c'}V_+;w_{P'})
\mathrm e^{-\vartheta_0d(P,P')}$ times
$\sum_{n\ge2}(\alpha_G\rho_G(V_+;w_\delta))^{n-2}\le2$.
This proves \eqref{eq:interacting-tree-mixed}.
\end{proof}

For a function $f$ of the marks, let
$t^T$ keep the coordinates indexed by $T$
and set all others to zero, and define
its component with variable set $S$ by
\[
 f_S(t)=\sum_{T\subset S}(-1)^{|S\setminus T|}f(t^T).
\]
Then $f=\sum_Sf_S$, $f_\varnothing=f(0)$,
and $f_S$ vanishes if any variable of $S$
vanishes. The latter statement follows by
pairing the terms with and without that
variable. In particular a variable absent
from $f$ makes every component containing
it zero. Two coordinate differences can
be represented by two integrals of the
mixed derivative. Consequently
\begin{equation}\label{eq:interacting-Mobius-bound}
 \sup_{|t|_\infty\le R}|f_S(t)|
 \le2^{|S|-2}R^2
          \sup_{|t|_\infty\le R}|\partial_c\partial_{c'}f(t)|
 \qquad(c,c'\in S,\ c\ne c').
\end{equation}

\begin{proposition}\label{prop:interacting-mixed-logarithm}
Assume the factor and mass conditions
\eqref{eq:interacting-factor-conditions}--\eqref{eq:interacting-small-complex-mass},
the tree conditions \eqref{eq:interacting-tree-smallness}, and
$d_{\rm sep}\ge6M+8$. For $|S|\ge2$, the normalized logarithm satisfies
\begin{equation}\label{eq:interacting-mixed-logarithm}
 \sup_{|t|_\infty\le t_x/2}|\mathcal G_S(t;\theta,A)|
                     \le\varpi_{\rm sep}\le\varpi_{\rm sep}^+,
\end{equation}
where
\begin{align}
 C_W^{\mathrm{mix}}&=3C_V^{\mathrm{ins}}h_G^2+2\alpha_G(C_V^{\mathrm{ins}})^2h_G^4
                             +4\exp(6)C_dM^3K_{\mathrm{pol}},\notag\\
 \varpi_{\rm sep}&=2^{N_*}C_W^{\mathrm{mix}}p_\beta
                    \mathrm e^{-\vartheta_0(d_{\rm sep}-4M-6)},\notag\\
 \varpi_{\rm sep}^+&=2^{N_*}C_W^{\mathrm{mix}}p_\beta
                    \mathrm e^{-\vartheta_0(Mr_*/(8L)-4M-7)} .
                                  \label{eq:interacting-separation-errors}
\end{align}
\end{proposition}

\begin{proof}
The tree representation and uniqueness of the
normalized logarithm give
\[
 \mathcal G(t)=\sum_X\mathfrak P_*(X,\mathscr C(t))+W(t)-W(0).
\]
The component of a determinant activity is zero
unless $X$ meets every $P_c$, $c\in S$.
Such a polymer has
$d_M(X)\ge d_{\rm sep}/M-2\ge1$.
Use \eqref{eq:interacting-P-bound}, sum the
polymers anchored at one of at most $p_\beta$
cubes, and keep $\mathrm e^{-2d_M(X)}$ from
$\kappa-\kappa_0\ge2$.
Its total contribution is at most
$2^{|S|}2\exp(4)C_dM^3K_{\mathrm{pol}}p_\beta
\mathrm e^{-\vartheta_0d_{\rm sep}}$.

Apply \eqref{eq:interacting-tree-mixed}
with $P=P_c$, $P'=P_{c'}$ and
\eqref{eq:interacting-differentiated-norms}.
The bound for $\partial_c\partial_{c'}W$ is
\[
 q_x^2p_\beta\left[
       2C_V^{\mathrm{ins}}h_G^2\mathrm e^{-\vartheta_0(d_{\rm sep}-4M-6)}
           +2\alpha_G(C_V^{\mathrm{ins}})^2h_G^4\mathrm e^{-\vartheta_0d_{\rm sep}}\right].
\]
In \eqref{eq:interacting-Mobius-bound}, take
$R=t_x/2$ and use $q_xR=1$.
Add the determinant contribution and
$|S|\le N_*$ to obtain the first inequality.
The placement estimate
$d_{\rm sep}\ge Mr_*/(8L)-1$
proves the second.
When the lattice probe is zero,
$R_{\rm tm}=1$, $\mathcal G=0$ and the
first separation error is interpreted as zero.
\end{proof}

\section{The positive small-field measure and its complement}
\label{subsec:interacting-complement}

\cite[Proposition~3.1]{TerminalPressure} identifies $F_N$
with $(A_*;\zeta_j)_{j<k_*}$. Their scalar
coordinates have the form
$\zeta_j(\Upsilon)=\langle\phi_{j,\Upsilon},F^{\rm u}\rangle$
with orthonormal $\phi_{j,\Upsilon}$; each
coarse plaquette coordinate $(dA_*)(p)$ has
the form $\langle\phi_p,F^{\rm u}\rangle$ with
$\|\phi_p\|\le1$.
There are at most $2\ell_*^3L^{3(k_*-j)}$
coordinates at level $j$ and $3\ell_*^3$
coarse plaquettes. The stop Gaussian
$\dd\mathsf m_*$ is nondegenerate on the
real mean-zero coclosed coarse fields.
The row and count statements are
\cite[(13)]{TerminalPressure} and
\cite[(14)]{TerminalPressure}; the
Gaussian assertion is \cite[(12)]{TerminalPressure}.

Let $\mathfrak G$ require
$|\zeta_j(\Upsilon)|\le p_{0,j}$
and $|dA_*|_\infty\le p_*$, and let
$\mathfrak G^-$ decrease each threshold by one.
Put
\begin{align}
 \varpi_0&=18\ell_*^3\mathrm e^{-p_{0*}/2}
               +6\ell_*^3\mathrm e^{-(p_*-1)^2/(2C_0^2)},\notag\\
 t_2&=e\,\min\left\{N_*^{-1/2},\,C_\beta^{-1},
                                     (2\sqrt{N_*\mathbf D})^{-1}\right\},\notag\\
 \varpi_L&=\sqrt{\varpi_0}\,\mathrm e^{\mathbf D/2},
 \qquad M_L=2+\mathrm e^{\mathbf D/2},\notag\\
 \lambda_L&=(\varpi_L^{1/3}M_L^{2/3})^{1/N_*}
                              5(2N_*r/t_2)^4 .
                                     \label{eq:interacting-complement-constants}
\end{align}
A reciprocal with zero denominator is interpreted
as $+\infty$. The number $\mathbf D$ is the
uniform probe energy bound, so $v_\eta\le\mathbf D$.

\begin{proposition}\label{prop:interacting-complement}
Assume the hypotheses of \cref{prop:interacting-one-cell,prop:interacting-mixed-logarithm}
and the tail conditions $v_*\ge8$ and $p_0v_*^{p_0-1}\ge28$.
For a fixed $0<s<1$, set
$\bar\kappa_1=6e_*^{1/4-9\varepsilon}$.
Assume
\begin{equation}\label{eq:interacting-jet-conditions}
 2\varpi_{\rm sep}^+\le\bar\kappa_1^{N_*},\quad
 N_*\bar\kappa_1^{1-s}\le1/2,\quad
 \varpi_0\le1/4,\quad 8N_*t_2\le t_x .
\end{equation}
There is a finite positive measure $v_{\rm sf}$
on $[-\pi,\pi]^3\times\{A\in V_{\mathrm t}(T_*^0):\chi_*(A)=1\}$, of
mass in $[1-\varpi_0,1]$, such that
\begin{equation}\label{eq:interacting-small-field-representation}
 \frac{\bar\nu^{a(t)}(\mathfrak G)}{\bar\nu(F_N)}
       =\int \mathrm e^{\mathcal G(t;\theta,A)}\,v_{\rm sf}(\dd\theta\,\dd A).
\end{equation}
The normalized complement
$L(t)=\bar\nu^{a(t)}(F_N\setminus\mathfrak G)/\bar\nu(F_N)$
is entire. Its fourth jet in each variable has
coefficient norm at radius $r$ at most
$\lambda_L^{N_*}$.
\end{proposition}

\begin{proof}
First, the union bound and
\cite[(7)]{TerminalPressure} give
\begin{equation}\label{eq:interacting-small-field-tail}
 \bar\mu(F_N\setminus\mathfrak G^-)\le\varpi_0 .
\end{equation}
For $i=k_*-j\ge1$, write
$v_j=v_*+i(\log L)/2$.
The scale conditions
$v_*\ge8$, $p_0v_*^{p_0-1}\ge28$ give
$p_{0,j}\ge p_{0*}+14i\log L$.
Since $(p_{0,j}-1)^2\ge p_{0,j}$,
the individual tail is at most
$2\mathrm e^{-p_{0*}/2}L^{-7i}$.
Multiply by the coordinate count
$2\ell_*^3L^{3i}$ and sum $L^{-4i}$.
The $3\ell_*^3$ coarse plaquettes give the
second term in \eqref{eq:interacting-complement-constants}.
These estimates yield \eqref{eq:interacting-small-field-tail},
with room in the coefficient $18$.

The exact density identity defines
\begin{equation}\label{eq:interacting-small-field-measure}
 v_{\rm sf}(\dd\theta\,\dd A)=
   \frac{\chi_*(A)Y_*^{\rm x}(A+c_\theta;0)}
             {\bar\nu(F_N)}\,\frac{\dd\theta}{(2\pi)^3}
                                      \dd\mathsf m_*(A).
\end{equation}
We check positivity of this density, which
does not follow from the formal Grassmann
representation. For real marks, the original
fine law $\nu^{\theta,a(t)}$ is nonnegative.
The identity \cite[(21)]{TerminalPressure} says, for every
nonnegative bounded Borel $\Phi$,
\[
 \int\Phi(A)\chi_*(A)Y_*^{\rm x}(A+c_\theta;\mathscr C(t))
                \,\dd\mathsf m_*(A)
       =\nu^{\theta,a(t)}(\mathbf1_{\mathfrak G}\Phi(A_*))\ge0.
\]
The terminal density is continuous. On the
interior, relative to $V_{\mathrm t}(T_*^0)$, of the compact
convex set $\{A\in V_{\mathrm t}(T_*^0):|dA|_\infty\le p_*\}$, a
nonzero imaginary part or negative real
part would persist on a ball of positive
Gaussian measure, contradicting this identity.
At the boundary, approach along $uA$,
$u\uparrow1$; zero is an interior point.
Thus the density is nonnegative everywhere
on this set. \Cref{prop:interacting-factor}
makes it nonzero there, hence strictly positive.
This reasoning also proves
$\mathcal G(t;\theta,A)\in\R$ for real $t$:
its exponential is a positive ratio and its
continuous value at zero is zero.

Integration of the exact terminal identity over
$\theta$, followed by
\eqref{eq:interacting-terminal-ratio}, proves
\eqref{eq:interacting-small-field-representation}.
At $t=0$ the mass is $\bar\mu(\mathfrak G)$,
which lies in $[1-\varpi_0,1]$.
All integrals here are well-defined Borel
integrals: the density is continuous on
the compact support and zero elsewhere.

For real $t=ez$ with $|t|_\infty\le t_2$,
put $G_z=\sum_cz_cd\widehat{\mathfrak a}_c$.
Disjoint supports imply
$\|G_z\|^2=v_\eta|z|_2^2\le1/4$.
The real source shift gives
\[
 L(ez)=\mathrm e^{-v_\eta|z|_2^2/2}
 \bar\mu\left(\mathbf1_{(F_N\setminus\mathfrak G)+G_z}
                              \mathrm e^{\langle G_z,F^{\rm u}\rangle}\right).
\]
Every photon coordinate changes by at most
$\|G_z\|\le1$ under the shift. Thus the
shifted complement is contained in
$F_N\setminus\mathfrak G^-$.
Cauchy--Schwarz, \eqref{eq:interacting-small-field-tail},
and Gaussian domination at $2G_z$ give
\begin{equation}\label{eq:interacting-real-complement}
 0\le L(t)\le\sqrt{\varpi_0}\,
          \mathrm e^{v_\eta|z|_2^2/2}\le\varpi_L .
\end{equation}
This is the only place where the lowered
thresholds of $\mathfrak G^-$ are needed.

The source-shift identity also gives, for complex marks,
\[
 |\mathring\Phi(t)|\le
       \exp\left(\frac{v_\eta}{2e^2}
                                      \sum_c(\operatorname{Im}t_c)^2\right).
\]
It follows directly by bounding the expectation
of the real exponential and canceling its
real quadratic normalization.
On $|t|_\infty\le t_2$ this is at most
$\mathrm e^{\mathbf D/2}$.
The small-field integral has modulus at most
$2$. Indeed, evenness and the double zero of
$\ell_0$ give $|\ell_0(t)|\le B_0|t|^2/t_x^2$,
while the zero of $w_c$ gives
$|w_c(t)|\le(3/8)|t|/t_x$.
Their sum over cells at
$t_2\le t_x/(8N_*)$ is at most
$1/(256N_*)+3/64$. The remaining logarithm
components sum to at most
$2^{N_*}\varpi_{\rm sep}^+
\le(2\bar\kappa_1)^{N_*}/2\le N_*^{-N_*}/2$,
by \eqref{eq:interacting-jet-conditions}.
Thus $|\mathcal G|\le1/2$ and
$|\mathrm e^{\mathcal G}|\le2$.
Since $L$ is the difference between the
full generating function and the small-field
integral, $|L|\le M_L$ on that polydisc.
It is entire because the finite determinant
density is entire and is bounded on compact
mark sets uniformly in the real photon
field and real twist. The modulus bound
in the external-response section supplies
the required Gaussian integrable majorant
also for either Borel event.

We use the following two-constants estimate.
Let $n\ge1$ be an integer. If a holomorphic function
on a polydisc in $\C^n$ of radius $R$ has modulus at most $M$, and
at most $\delta\le M$ on its real points,
then its modulus is at most $\delta^{1/3}M^{2/3}$
on the polydisc of radius $R/(2n)$.
On the upper
half of the unit disc the harmonic function
$\omega(z)=\frac2\pi\arg[(1+z)/(1-z)]$
is zero on the diameter and one on the
semicircle. Comparing $\log|f(R'z)|$
with $(1-\omega)\log\delta+\omega\log M$,
for $R'<R$, the maximum principle gives
the same inequality in the half disc.
To handle its endpoints, subtract
$\epsilon\log[4/|1-z^2|]$ before applying
the principle, then let $\epsilon\downarrow0$.
Its harmonic conjugate gives a holomorphic
function to which the modulus principle
applies. On $|z|\le q$, $\omega(z)\le
(4/\pi)\arctan q$; reflection treats the
lower half disc. Let $R'\uparrow R$.
Apply this one-variable estimate successively
with $q=1/(2n)$. The remaining exponent of
$\delta$ is at least
$(1-\frac4\pi\arctan(1/(2n)))^n
\ge1-2/\pi>1/3$, proving the assertion.
The case $\delta=0$ follows by the identity
theorem, and $M=0$ is immediate.

Apply it to $L$, with $n=N_*$ and $R=t_2$.
Cauchy's coefficient bound gives
\[
 |L_\alpha|\le\varpi_L^{1/3}M_L^{2/3}(2N_*/t_2)^{|\alpha|}.
\]
In the fourth jet, every exponent is at most
four. Since $y=2N_*r/t_2\ge1$,
\[
 \sum_{\alpha_c\le4}|L_\alpha|r^{|\alpha|}
   \le\varpi_L^{1/3}M_L^{2/3}
                 (1+y+\cdots+y^4)^{N_*}
   \le\lambda_L^{N_*}.
\]
Reality on real marks gives real jet coefficients.
\end{proof}

\section{The fourth-order cancellation}
\label{subsec:interacting-fourth-jet}

Work in the finite ring
$\mathcal J=\C[t_c:c\in\mathcal C_*]/(t_c^5:c\in\mathcal C_*)$,
with coefficient norm
$\|P\|_r=\sum_{\alpha_c\le4}|P_\alpha|r^{|\alpha|}$.
This norm is submultiplicative, since discarding
terms with an exponent at least five can only
decrease the coefficient sum.
The jet map is a ring homomorphism; it
keeps degree at most four in each variable.
For $S=\{c\}$ set $k_S=\mathrm e^{w_c}-1$ in this
ring, and for $|S|\ge2$ set
$k_S=\mathrm e^{\mathcal G_S}-1$.
Every monomial of $k_S$ uses every variable
of $S$ and no other variable.

For a component $f_S$ bounded on a polydisc
of radius $R$, Cauchy's estimate at a smaller
radius $R'\le R/2$ gives
\[
 \|f_S\|_{R'}\le\sup|f_S|
             \left(\sum_{j=1}^4(R'/R)^j\right)^{|S|}
                         \le\sup|f_S|.
\]
Its monomials contain each variable of $S$
because it vanishes on every corresponding
coordinate hyperplane. Use this with
$R=t_x/2$, $R'=r$, and use the one-cell
estimate on radius $t_x$, to obtain
\begin{equation}\label{eq:interacting-error-jets}
 \|k_{\{c\}}\|_r\le\bar\kappa_1,\qquad
 \|k_S\|_r\le2\varpi_{\rm sep}^+
                       \le\bar\kappa_1^{N_*}\quad(|S|\ge2).
\end{equation}
We used $\|w_c\|_r\le3e_*^{1/4-9\varepsilon}\le3/8$
and $\|\mathrm e^P-1\|_r\le2\|P\|_r$ for $\|P\|_r\le1/2$.
The coefficients are continuous in $(\theta,A)$:
express the logarithm coefficients by Cauchy
integrals on a fixed small torus, and then
take the finite polynomial expressions in
those coefficients. They are real because
the logarithm is real at real marks.

In $\mathcal J$ the normalized generating
function has the exact representation
\begin{equation}\label{eq:interacting-jet-representation}
 \mathring\Phi=
   \prod_c\phi_0(t_c)
       \int\prod_{\varnothing\ne S\subset\mathcal C_*}
                         (1+k_S)\,v_{\rm sf}(\dd\theta\,\dd A)
       +L^{[4]}.
\end{equation}
To justify taking the jet inside the integral,
use joint continuity, holomorphy, and
$|\mathrm e^{\mathcal G}|\le2$ on a fixed
polydisc of radius at most $t_2$.
Cauchy's coefficient formulas and Fubini
then commute every required coefficient
with the finite measure integral.
Decompose $\mathcal G$ into its variable-set
components; its singleton components are
$\ell_0(t_c)+w_c(t_c)$.
Exponentiation, the jet homomorphism,
and $\mathrm e^{\ell_0}=N_0$ prove the representation.

\begin{theorem}\label{thm:interacting-reference-error}
Use the fixed construction, common trajectory and terminal inputs of
\cref{sec:rrm-terminal-inputs}. Suppose the factor and mass conditions
\eqref{eq:interacting-factor-conditions}--\eqref{eq:interacting-small-complex-mass},
the one-cell conditions \eqref{eq:interacting-one-cell-conditions},
the tree conditions \eqref{eq:interacting-tree-smallness}, and the jet
conditions \eqref{eq:interacting-jet-conditions} hold. Assume also
$r_*>31L+16$, $Mr_*/(8L)-1\ge6M+8$,
$v_*\ge8$, $p_0v_*^{p_0-1}\ge28$, and
$\lambda_L\le\bar\kappa_1^s$, where $0<s<1$.
For $m=2,4$,
\begin{equation}\label{eq:interacting-Theta-bound}
 \Theta_m^{(0)}\le
        m!r^{-m}\bigl(2\exp(1)\bar\kappa_1^s+\lambda_L\bigr)
                       \le8m!r^{-m}\bar\kappa_1^s.
\end{equation}
With
$\mathfrak e=s(1/4-9\varepsilon)-12\varepsilon>0$,
\begin{align}
 \Theta_4^{(0)}
    &\le3\cdot10^5 C_b^4e_*^{\mathfrak e}
           \le3\cdot10^5C_b^4h^{1/8}e^{\mathfrak e},\notag\\
 (1+|c_2|)\Theta_2^{(0)}
    &\le3\cdot10^4C_b^4e_*^{\mathfrak e}
           \le3\cdot10^4C_b^4h^{1/8}e^{\mathfrak e}.
                                    \label{eq:interacting-final-terminal-error}
\end{align}
The constants are independent of the ultraviolet cutoff and the
admitted placement. The terminal estimates used before harmonic
averaging hold uniformly in the real twist.
\end{theorem}

\begin{proof}
On the one-variable jet space define
$\ell_c(t_c^j)=m!\mathbf1_{\{j=m\}}-c_m\mathbf1_{\{j=0\}}$
and let $\ell=\bigotimes_c\ell_c$.
Then $(\Theta_m^{(0)})^{N_*}=|\ell(\mathring\Phi)|$.
The free jet satisfies $\ell_c(\phi_0)=0$.
From \eqref{eq:interacting-free-jet},
$|\ell(P)|\le(m!r^{-m})^{N_*}\|P\|_r$.
These statements follow by evaluating the
functionals on monomials with every exponent
in $\{0,m\}$.

Expand the product in
\eqref{eq:interacting-jet-representation} over
families $\mathcal Y$ of nonempty cell sets.
If $\mathcal Y$ fails to cover a cell $c$,
its error factors do not contain $t_c$;
the factor $\phi_0(t_c)$ is therefore
annihilated by $\ell_c$. Only covering
families contribute.
Assign cost one to each singleton and
cost $N_*$ to each larger set, and let
$a(\mathcal Y)$ be the sum of these costs.
A covering family has $a(\mathcal Y)\ge N_*$:
either it contains a larger set or it
contains every singleton.
By \eqref{eq:interacting-error-jets} its
coefficient norm is at most
$\bar\kappa_1^{a(\mathcal Y)}$.
The free factors have norm at most $2^{N_*}$.
Keep $\bar\kappa_1^{sN_*}$ from the cost
and sum the rest over all families. This
sum is
\[
 (1+\bar\kappa_1^{1-s})^{N_*}
  (1+\bar\kappa_1^{(1-s)N_*})^{2^{N_*}-1-N_*}
 \le\exp\left(N_*\bar\kappa_1^{1-s}
                  +2^{N_*}\bar\kappa_1^{(1-s)N_*}\right)
 \le\exp(1).
\]
The last step uses $N_*\bar\kappa_1^{1-s}\le1/2$.
The positive small-field measure has mass at
most one. Its total contribution is thus
bounded by
$(m!r^{-m})^{N_*}\exp(1)(2\bar\kappa_1^s)^{N_*}$.
The complement contributes at most
$(m!r^{-m}\lambda_L)^{N_*}$.
Taking the $N_*$th root proves
\eqref{eq:interacting-Theta-bound}.

Since $r^{-4}=256C_b^4e_*^{-12\varepsilon}$
and $\bar\kappa_1^s\le6e_*^{s(1/4-9\varepsilon)}$,
the $m=4$ bound has coefficient at most
$8\cdot24\cdot1536<3\cdot10^5$.
For $m=2$, $1+|c_2|\le(17/16)r^{-2}$,
so its coefficient is at most
$17\cdot1536<3\cdot10^4$.
Finally $0<\mathfrak e<1/4$ and
$e_*\le e\sqrt h$ imply
$e_*^{\mathfrak e}\le h^{1/8}e^{\mathfrak e}$.
The scale and probe bounds entering this
proof are uniform in the stated choices.
\end{proof}

\section{A positive charge interval}
\label{subsec:interacting-positive-threshold}

We now verify that the response conditions admit
a common interval of charges, independent of the
ultraviolet cutoff. In addition to
\eqref{eq:interacting-factor-conditions},
\eqref{eq:interacting-small-complex-mass},
\eqref{eq:interacting-one-cell-conditions},
\eqref{eq:interacting-tree-smallness}, and
\eqref{eq:interacting-jet-conditions}, require
$r_*>31L+16$, $v_*\ge8$, $p_0v_*^{p_0-1}\ge28$, and
\begin{gather}
 Mr_*/(8L)-1\ge6M+8,\qquad
 \lambda_L\le\bar\kappa_1^s,
                                  \label{eq:interacting-last-response-conditions}\\
 D_\beta^{\rm buf}/2+1/2\le3\ell_c/8,\qquad
 D_\beta^{\rm buf}+3\le\ell_c .
                                  \label{eq:interacting-placement-conditions}
\end{gather}
The construction of \cite[\S6.1]{CoverTuning}
fixes an odd $L\ge13$, a power
$M$ of $L$ with $M\ge L^2$, a localization
power $r_{\rm loc}\ge2$, and cutoff powers
$p_0\ge2$, $p>2p_0$. It also fixes
$0<\varepsilon\le1/200$ and $0<s<1$ such that
\begin{equation}\label{eq:interacting-response-exponents}
 \mathfrak e=s(1/4-9\varepsilon)-12\varepsilon>0.
\end{equation}
For the construction fixed in the introduction, take $s=1/2$;
then $\mathfrak e\ge17/400$.

\begin{proposition}\label{prop:interacting-positive-threshold}
For the fixed construction constants, $\bar m$,
and probe $\beta$, there is $e_\sharp(\beta)>0$
such that all response conditions hold for
$0<e\le e_\sharp(\beta)$ and every
$N\ge N_{\rm ref}(e)$. The threshold depends
on the probe only through
$C_\beta,\mathbf D,D_\beta^{\rm buf}$.
It is independent of $N$, the twist, and the
admitted placement. All centered placements
are admitted in this range.
\end{proposition}

\begin{proof}
Write
$v_K=-\log(e\sqrt h)$ and
$r_K=\min\{L^i:L^i\ge v_K^{r_{\rm loc}}\}$.
Then $R^*(e)=hMr_K/L$ and
$\ell_c=R^*(e)/(2L)$.
\cite[Proposition~6.1]{CoverTuning} and
\cite[(53)]{CoverTuning} give
a finite $N_{\rm ref}(e)$.

For $N\ge N_{\rm ref}(e)$ the stop level
is $K-1$ or $K-2$ by
\cite[(54)]{CoverTuning}. In particular its
scales are independent of $N$ and satisfy
\begin{gather}
 v_K\le v_*,\quad e_*=\mathrm e^{-v_*},\quad
 a_*\in\{h/L,h/L^2\},\notag\\
 v_*^{r_{\rm loc}}\le r_*,\quad
 Mv_*^{r_{\rm loc}}\le\ell_*<LMv_*^{r_{\rm loc}},
 \quad p_{0*}=v_*^{p_0},\quad p_*=v_*^p,\notag\\
 p_\beta\le p_\beta^+:=
                (D_\beta^{\rm buf}L^2/h+2M+4)^3,\qquad
 \ell_c\ge hMv_K^{r_{\rm loc}}/(2L^2).
                                  \label{eq:interacting-threshold-scales}
\end{gather}
The strict wrapping condition $r_*>31L+16$
already follows from \cite[(50)]{CoverTuning}
whenever $e\sqrt h\le e_1$.
Thus every response condition is a condition
on $e$ alone once the construction and probe
are fixed.

A bound $A v^d \mathrm e^{-\gamma v}$
with $A,d\ge0$ and $\gamma>0$ tends to zero,
uniformly on a sufficiently large half-line.
Split the exponential in half and use
$v^d\mathrm e^{-\gamma v/2}\le(2d/(\gamma\exp(1)))^d$.
Likewise, $cv^2$ exceeds $bv+\log^+A$
on a sufficiently large half-line for
$c>0$. We apply these two inequalities to
the following groups of conditions.

The local tree smallness and terminal factor
conditions follow from exponential decay
with positive rates
$\varepsilon$, $2\varepsilon$,
$1/4-\varepsilon$, and $1/4-9\varepsilon$;
the volume factors are bounded by
$(LM)^3v_*^{3r_{\rm loc}}$.
For $\bar r$, use
$1/4-7\varepsilon>2\varepsilon$.
The complex mass and coarse determinant
conditions similarly have rate $2\varepsilon$.
The support count is bounded by $p_\beta^+$.
All three summands of $B_0$ tend to zero:
their rates are $4\varepsilon$, $2\varepsilon$,
and the quadratic decay $\mathrm e^{-r_*}\le \mathrm e^{-v_*^2}$.
Also $\tau_s\ge c_\tau Mv_*^2$, so both
$\tau_s\ge1$ and
$52B_1\mathrm e^{-\tau_s}\le e_*^{1/4-9\varepsilon}$
hold once $v_*$ is large.

The separation condition follows from
$r_*\ge v_*^2$. Its error obeys
\[
 2\varpi_{\rm sep}^+
 \le2^{N_*+1}C_W^{\mathrm{mix}}p_\beta^+
          \mathrm e^{\vartheta_0(4M+7)}
          \mathrm e^{-\vartheta_0Mv_*^2/(8L)} .
\]
Since $\bar\kappa_1\ge \mathrm e^{-v_*/4}$,
quadratic decay makes this at most
$\bar\kappa_1^{N_*}$ for all sufficiently
large $v_*$. The family summation condition
$N_*\bar\kappa_1^{1-s}\le1/2$
uses the positive rate
$(1-s)(1/4-9\varepsilon)$.

For the photon complement, the inequalities $v_*\ge8$ and
$p_0v_*^{p_0-1}\ge28$ hold on a sufficiently large half-line,
since $p_0\ge2$. Moreover $p>2p_0$ implies
$(p_*-1)^2\ge C_0^2p_{0*}$ once $v_*$
is large. Therefore
\begin{equation}\label{eq:interacting-complement-threshold}
 \varpi_0\le24(LM)^3v_*^{3r_{\rm loc}}\mathrm e^{-v_*^2/2}.
\end{equation}
This is eventually at most $1/4$.
Moreover $t_2/e$ is fixed and positive, while
$e$ is between $(L/h)^{1/2}e_*$
and $Lh^{-1/2}e_*$. Hence
$8N_*t_2/t_x$ is bounded by a fixed
constant times $e_*^{1-3\varepsilon}$,
which tends to zero.

The potentially largest loss is the fourth
derivative radius in $\lambda_L$.
From \eqref{eq:interacting-complement-threshold},
the definition of $\lambda_L$, and
$r/t_2\le C \mathrm e^{v_*}$, we obtain
\begin{equation}\label{eq:interacting-lambda-threshold}
 \lambda_L\le A_8
   \exp\left(-\frac{v_*^2}{12N_*}
                     +\left(\frac{r_{\rm loc}}{2N_*}+4\right)v_*\right)
\end{equation}
for a finite $A_8>0$ depending only on the
fixed constants and probe bounds.
For clarity one can take
\[
 A_8=[24(LM)^3]^{1/(6N_*)}
       \mathrm e^{\mathbf D/(6N_*)}M_L^{2/(3N_*)}
       5\left[
        \frac{N_*\sqrt h}{2C_b\sqrt L}(t_2/e)^{-1}\right]^4 .
\]
If
$v_*^2/(12N_*)\ge
(r_{\rm loc}/(2N_*)+17/4)v_*+\log^+A_8$,
then $\lambda_L\le \mathrm e^{-v_*/4}\le\bar\kappa_1^s$.
Thus the complex continuation and the
fourth derivatives preserve a positive
small-charge interval.
Finally the two placement inequalities
hold when
$\ell_c\ge(4/3)D_\beta^{\rm buf}+3$;
this follows from
\eqref{eq:interacting-threshold-scales} as $v_K\to\infty$.

Choose a finite $V$ beyond all the half-lines
just specified and put
$e_V=h^{-1/2}\min(e_1,\mathrm e^{-V})>0$.
Every $0<e\le e_V$ has $v_*\ge v_K\ge V$
and satisfies every response condition.
Define $e_\sharp$ as one half of the
supremum of the charges $e'\le e_1h^{-1/2}$
for which all conditions hold for every
$0<e\le e'$. Its defining set contains
$e_V$, so $e_\sharp\ge e_V/2>0$.
For $e\le e_\sharp$, a member of the
defining set exceeds $e_\sharp$, which
ensures that all inequalities hold at $e$
even if the scale functions have jumps.
This avoids assuming monotonicity of the
rounded stop scales.

The scale identities depend on $e,h$ and
the fixed construction constants; the probe
enters through
$C_\beta,\mathbf D,D_\beta^{\rm buf}$.
They do not involve the twist or placement,
and they are unchanged by the coordinate
reflections used to form the cell probes.
Centered placements satisfy the admission
inequality by
\eqref{eq:interacting-placement-conditions},
since their centers differ from the cell
center by at most $\eta/2\le1/2$.
The same conditions give
$D_\beta^{\rm buf}+3\eta\le\ell_c$ and
$R^*(e)\ge D_\beta+1$.
\end{proof}



\begin{proof}[Proof of \cref{thm:rrm-main}]
The base construction gives the positive physical ceiling
\[
 e_G=\begin{cases}
 e_1,&h=1,\\
 \min\{e_1/\sqrt h,\,1/(\sqrt h\log h)\},&h>1,
 \end{cases}
\]
as in \cite[(31)]{QED3Integration}. On this interval the common pair
and all upstream terminal estimates apply.
Choose $e_{\rm RR}=\min\{e_G,e_\sharp(\beta)\}$, with the threshold
of \cref{prop:interacting-positive-threshold}, at $s=1/2$.
That proposition verifies the additional inequalities of
\cref{thm:interacting-reference-error} for every smaller positive
charge and every $N\ge N_{\rm ref}(e)$, and admits every centered
placement. The bounds \eqref{eq:interacting-final-terminal-error}
give \eqref{eq:rrm-main-bound} with the stated common constant.
The exact unfolding identifies the normalized density with the
zero-flux law, \cref{prop:external-entire} proves entireness, and
\eqref{eq:external-Wick-product} proves \eqref{eq:rrm-main-wick}.
Its differentiation proof establishes the required integrability.
\end{proof}
\printbibliography
\end{document}
